\chapter{语言任务的求解与执行}
\label{chap:nlp-solving-and-execution}

“这次读书会改到什么时候？”需要依据通知回答；“把我的预约改到通知中的时间”却要求另一种交付物：预约记录真的发生了变化，而且没有改错对象。第\ref{chap:nlp-question-answering-dialogue}章讨论了怎样取得证据、澄清要求和维护对话状态。本章处理输出须被执行或检查的任务：数学与逻辑题的解答要能代回原题复算或由证明助手检查，程序要在约定输入上表现出规定的行为，环境中的行动要使外部状态真正达到目标。三类任务都在“生成候选—执行—依据反馈修订”的循环中完成，但验收依据不同：算术题核对表达式、单位和原始条件，程序任务核对测试与规格，环境任务读取行动后的实际状态。

这类任务的困难可以归纳为八项挑战。
\begin{enumerate}
  \item 把语言条件准确转成形式对象：数量、折扣作用域、“恰好”与“至少”之差，一旦映射错误，后面再可靠的计算也只是忠实地算错。
  \item 让计算交给可靠的执行器：语言模型逐词生成的算术在长计算中容易出错，解释器、求解器与证明助手则按确定的规则执行。
  \item 验证答案和中间步骤：答案正确不代表过程正确，自然语言证明又无法用数值回代。
  \item 生成语法和语义都正确的代码：单次采样的正确率有限，语法错误可以由约束排除，语义错误却只能借执行与测试发现。
  \item 在大型仓库中定位修改位置：问题描述与需要修改的代码之间隔着大量文件和调用关系。
  \item 获得调用工具的能力并填对参数：模型要知道有哪些工具、何时调用、每个参数的值从哪里来。
  \item 在多步交互中规划与记忆：行动顺序要随观察调整，上下文有限时还要保留仍影响任务的事实。
  \item 从失败中恢复并判断完成：参数错误、超时与版本冲突需要不同的处理，“执行器说完成了”与“环境已经达到目标”也不是一回事。
\end{enumerate}

应对这些挑战的方法大致沿一条主线发展：从自然语言逐步推理，到生成可执行的中间表示（表达式、约束、程序、形式证明），再到依据执行反馈修订和搜索，最后是在环境中交替判断、行动和观察的智能体及其训练。每一步都把一部分原本由模型“凭语言能力”承担的工作交给可检查的外部机制：中间表示让计算可以由执行器完成，执行反馈让错误在交付前暴露，环境观察让行动依据真实状态而不是模型的设想。代价也随之出现：形式化本身可能出错，执行与搜索消耗调用次数和时间，环境中的写操作还可能产生无法撤回的副作用。

本章与前文的分工如下。思维链、自一致性、思维树与蒙特卡洛树搜索等通用的推断时推理方法，以及结果奖励模型、过程奖励模型和基于可验证奖励的强化学习，在第\ref{chap:nlp-reasoning}章主讲；pass@$k$的定义与无偏估计在第\ref{chap:nlp-evaluation}章；约束解码的有限状态机与下推自动机机制在第\ref{sec:nlp-inference-grammar}小节；函数调用的对话格式在第\ref{sec:nlp-dialogue-llm-function}小节。本章回指这些内容，只展开它们接入具体任务时需要补充的结构、数据与算法。

章内顺序按“交付物离外部世界有多远”排列。第\ref{sec:nlp-execution-solving}节的数学与逻辑题输入固定、不改变外部状态，最适合先建立“形式化—执行—核验”这组机制，并引入程序执行与形式证明两种可靠的检查手段。第\ref{sec:nlp-execution-programs}节把程序从求解的中间工具变为交付物本身，执行反馈从“算出一个数”变成“通过一组测试”，修改对象也从一段代码扩大到整个仓库。第\ref{sec:nlp-execution-environment}节进入工具与环境，行动会改变外部状态，规划、记忆、失败恢复和完成判断成为新的问题。文中的数量题、程序仓库、工具接口和运行轨迹均为教学构造；所引论文提供方法依据，不为这些构造案例的效果背书。

\section{数学与逻辑问题求解}
\label{sec:nlp-execution-solving}

“算一个金额”“找一组满足排班规则的取值”和“证明一个命题”交付的是三种不同的产物：一个数、一组赋值、一个证明。它们的成功判断也各不相同：金额可以代回原题复算，赋值可以逐条检查约束，证明却无法用数值核对，只能由检查器验证推理的每一步。这类任务的输入是固定的题目，求解不改变外部状态，候选可以反复构造和检查，因此适合先研究“生成—执行—核验”这一组机制。

本节要处理四项挑战。第一，语言条件可能被误读：折扣施加给了错误的对象，或“恰好一人”被读成“至少一人”。第二，即使关系写对，长计算中的算术也会出错。第三，答案正确不代表过程正确，碰巧得到的答案无法支持解释。第四，证明不能用数值核对，自然语言写成的证明又难以逐步检查。方法谱系沿这四项挑战展开：把题目转成表达式树或约束，交给确定的算法或求解器（第\ref{sec:nlp-execution-formalization}小节）；让模型写程序并由解释器执行，再让文本推理与代码执行交替进行（第\ref{sec:nlp-execution-program-aided}小节）；用回代、规则和奖励模型检查答案与过程（第\ref{sec:nlp-execution-verification}小节）；把命题形式化后在证明助手中搜索并检查证明（第\ref{sec:nlp-execution-theorem-proving}小节）。节末表~\ref{tab:nlp-execution-solving-methods}对照各方法解决的问题与新增的代价。

\subsection{从语言到形式对象}
\label{sec:nlp-execution-formalization}

语言中的量、对象和条件需要成为计算能够读取的输入。第\ref{chap:nlp-text-analysis}章讨论的数量识别、指代、语义关系和时间规范化在这里提供映射依据，本章还要记录每个形式对象来自哪一项条件。这样就能区分三类失败：缺少单价属于信息不足；把“不打折”读成“打折”属于解析错误；在已正确建立的约束上没有找到解，才属于求解问题。形式对象一旦建立，计算就可以交给确定的算法，错误便集中到形式化这一步。下面依次讨论表达式树、有限域约束和模理论约束三种对象，它们能表达的问题依次扩大：表达式树只描述一条计算，有限域约束描述一组离散取值须同时满足的条件，模理论约束再加入整数、实数等数域上的运算与逻辑组合。

\subsubsection{表达式树与目标驱动树解码}
\label{sec:nlp-execution-expression}

考虑一道数量题：“购买3本书，每本25元；另买5支笔，每支4元。书打八折，笔不打折，不收其他费用。应付多少元？”一个可检查的表示应同时保存数值、单位和修饰范围。设书的单价为$p_{\mathrm{b}}=25$元/本、数量为$n_{\mathrm{b}}=3$本，笔的单价为$p_{\mathrm{p}}=4$元/支、数量为$n_{\mathrm{p}}=5$支，书的支付比例为$r_{\mathrm{b}}=0.8$。比例没有单位，“八折”指支付原价的80\%，不是减去原价的80\%。

\term{表达式树}（expression tree）把已知量作为叶，把运算作为内部节点，根节点给出待求的量。这里先分别计算书与笔的金额，只对书的分支乘折扣，再相加：
\begin{equation}
  a=(p_{\mathrm{b}}n_{\mathrm{b}})r_{\mathrm{b}}
      +p_{\mathrm{p}}n_{\mathrm{p}}
    =(25\times3)\times0.8+4\times5=80\text{元}.
  \label{eq:nlp-execution-price}
\end{equation}
这棵树可写为$\operatorname{Add}(\operatorname{Mul}(\operatorname{Mul}(25,3),0.8),\operatorname{Mul}(4,5))$。树一旦确定，求值只是自底向上的四则运算；难点在于从题目文字得到这棵树。类型检查能排除$25\text{元/本}+3\text{本}$这类单位不相容的组合，却排除不了$(25\times3+4\times5)\times0.8=76$元：后一式单位完全正确，只是把折扣错误地施加给了笔。若将“另买5支”错绑定为“买3支”，还会得到72元。因此，恢复树以后还要逐项核对数量指代、折扣作用域和“无其他费用”这项封闭条件。

\paragraph{序列到序列求解}
最早的神经求解器把表达式当作词元序列来生成。Wang等的深度神经求解器先把题目中的数字替换为按出现顺序编号的占位词元$n_0,n_1,\ldots$，再用循环神经网络编码器—解码器把题目翻译成含占位词元的表达式，最后代回数值计算\scite{nlp-execution-wang2017dns}。本例中$n_0=3,n_1=25,n_2=5,n_3=4$；“八折”在数量规范化时转成$n_4=0.8$。目标序列可以是中缀式“$(n_1\times n_0)\times n_4+n_3\times n_2$”。这种做法不需要人工特征，能生成训练集中没有出现过的表达式，但有两个结构上的缺陷。其一，解码器逐个输出括号和符号，必须自己学会括号配对与运算元个数，可能生成“$((n_1\times n_0)\times n_4$”这样无法求值的序列。其二，解码器只看到线性前缀，不知道当前词元在树中处于哪个位置、在为哪个子量服务；同一棵树的不同线性化（中缀、后缀、交换律下的不同写法）又会给出互相冲突的监督。

\paragraph{目标驱动树解码}
\term{目标驱动树结构模型}（goal-driven tree-structured model, GTS）改为直接生成表达式树\scite{nlp-execution-xie2019gts}。它的出发点是人解题的方式：先明确要求的量是什么，再判断能否直接从已知数量得到；不能时，就把它拆成由一个运算符连接的两个子量，对每个子量重复这一过程，直到每个子量都能由题中的某个数量或常数直接给出。每个树节点因此关联三样东西：一个\term{目标向量}（goal vector）$\bm{q}\in\mathbb{R}^{d}$，表示“这棵子树要算出什么”；一个预测的词元$\hat y$；以及子树完成后的\term{子树嵌入}（subtree embedding）$\bm{t}\in\mathbb{R}^{d}$，表示“这棵子树实际算的是什么”。目标自顶向下分解，子树嵌入自底向上合成，信息在树上双向流动。

先看输入与候选词元。设题目有$n$个词，编码器是双向门控循环单元（GRU），两个方向的隐藏状态相加得到$\bm{h}_s\in\mathbb{R}^{d}$，$s=1,\ldots,n$，原文取$d=512$。根目标$\bm{q}_0$取前向最后一个状态与后向第一个状态之和。解码词表因题而异：$V^{\mathrm{dec}}=V_{\mathrm{op}}\cup V_{\mathrm{con}}\cup N_P$，其中$V_{\mathrm{op}}$是运算符，$V_{\mathrm{con}}$是题中不出现但解答可能用到的常数（如1、$\pi$），$N_P$是本题的数量占位词元。词元$y$的嵌入$\bm{e}(y)\in\mathbb{R}^{d}$按类别取得：运算符与常数查各自的可训练嵌入矩阵，数量$n_k$则直接取它在题目中位置上的编码器状态。数量嵌入随题目上下文变化，模型“指向”题中的某个数而不是从固定词表中选数，这与式\eqref{eq:nlp-analysis-copy}的复制机制作用相同。

在每个节点，模型用目标向量对题目做注意力，得到上下文向量$\bm{c}=\sum_s\alpha_s\bm{h}_s$，其中$\alpha_s\propto\exp\bigl(\bm{v}_{\mathrm{a}}^{\top}\tanh(\bm{W}_{\mathrm{a}}[\bm{q};\bm{h}_s])\bigr)$。随后对每个候选词元打分并归一化：
\begin{equation}
  \begin{aligned}
  s(y\mid\bm{q},\bm{c})&=\bm{w}_{\mathrm{n}}^{\top}\tanh\bigl(\bm{W}_{\mathrm{s}}[\bm{q};\bm{c};\bm{e}(y)]\bigr),\\
  p(y\mid\bm{q},\bm{c})&=\frac{\exp s(y\mid\bm{q},\bm{c})}{\sum_{y'\in V^{\mathrm{dec}}}\exp s(y'\mid\bm{q},\bm{c})},
  \end{aligned}
  \label{eq:nlp-execution-gts-score}
\end{equation}
其中$[\cdot;\cdot]$表示拼接，$\bm{W}_{\mathrm{s}}\in\mathbb{R}^{d\times3d}$，$\bm{w}_{\mathrm{n}}\in\mathbb{R}^{d}$。打分把“当前目标”“题中相关信息”和“候选词元”放在一起比较；由于数量词元的嵌入来自题目，同一套参数可以对不同题目的不同数量打分。预测出的词元决定这个目标怎样实现：若是数量或常数，目标由它直接实现，该节点成为叶；若是运算符，目标要分解为左、右两个子目标。

左子目标只依赖父节点的信息。原文用两层带门控的前馈网络：
\begin{equation}
  \begin{aligned}
  \bm{u}&=[\bm{q};\bm{c};\bm{e}(\hat y)],\\
  \bm{h}_{\mathrm{l}}&=\sigma\bigl(\bm{W}_{\mathrm{ol}}\bm{u}\bigr)\odot\tanh\bigl(\bm{W}_{\mathrm{cl}}\bm{u}\bigr),\\
  \bm{q}_{\mathrm{l}}&=\sigma\bigl(\bm{W}_{\mathrm{gl}}\bm{h}_{\mathrm{l}}\bigr)\odot\tanh\bigl(\bm{W}_{\mathrm{le}}\bm{h}_{\mathrm{l}}\bigr),
  \end{aligned}
  \label{eq:nlp-execution-gts-left}
\end{equation}
其中$\sigma$是逐元素的sigmoid，$\odot$是逐元素乘积，$\bm{W}_{\mathrm{ol}},\bm{W}_{\mathrm{cl}}\in\mathbb{R}^{d\times3d}$，$\bm{W}_{\mathrm{gl}},\bm{W}_{\mathrm{le}}\in\mathbb{R}^{d\times d}$。sigmoid一支决定保留父信息的哪些分量，tanh一支给出候选内容，两者相乘便是门控。

右子目标要等左子树完全生成以后才计算，并额外读取左子树的嵌入$\bm{t}_{\mathrm{l}}$：
\begin{equation}
  \begin{aligned}
  \bm{h}_{\mathrm{r}}&=\sigma\bigl(\bm{W}_{\mathrm{or}}\bm{u}\bigr)\odot\tanh\bigl(\bm{W}_{\mathrm{cr}}\bm{u}\bigr),\\
  \bm{q}_{\mathrm{r}}&=\sigma\bigl(\bm{W}_{\mathrm{gr}}[\bm{h}_{\mathrm{r}};\bm{t}_{\mathrm{l}}]\bigr)\odot\tanh\bigl(\bm{W}_{\mathrm{re}}[\bm{h}_{\mathrm{r}};\bm{t}_{\mathrm{l}}]\bigr),
  \end{aligned}
  \label{eq:nlp-execution-gts-right}
\end{equation}
这里$\bm{W}_{\mathrm{gr}},\bm{W}_{\mathrm{re}}\in\mathbb{R}^{d\times2d}$。为什么右目标要看左子树？以加法为例，父目标“应付总额”分解出的两个子目标在父节点看来是对称的，若只用父信息生成，右目标可能与左目标几乎相同，模型会在两边生成同一个子量。读取$\bm{t}_{\mathrm{l}}$以后，右目标知道左边已经算出了“书的折后金额”，剩下的应该是“笔的金额”。

子树嵌入用自底向上的递归组合计算：叶节点的嵌入就是其词元嵌入$\bm{e}(\hat y)$；运算符节点的左右子树都完成后，
\begin{equation}
  \bm{t}=\sigma\bigl(\bm{W}_{\mathrm{gt}}[\bm{t}_{\mathrm{l}};\bm{t}_{\mathrm{r}};\bm{e}(\hat y)]\bigr)\odot
  \tanh\bigl(\bm{W}_{\mathrm{ct}}[\bm{t}_{\mathrm{l}};\bm{t}_{\mathrm{r}};\bm{e}(\hat y)]\bigr),
  \label{eq:nlp-execution-gts-comb}
\end{equation}
其中$\bm{W}_{\mathrm{gt}},\bm{W}_{\mathrm{ct}}\in\mathbb{R}^{d\times3d}$。目标向量描述的是“打算算什么”，子树嵌入描述的是“已经算出了什么”；后者由实际生成的运算符与数量合成，比目标向量更贴近已完成的计算。

生成顺序是树的前序遍历：先生成节点本身，再递归生成左子树，最后生成右子树。训练时使用标注的表达式树做教师强制，模型在前序的每个位置预测正确词元，目标函数是整棵树的负对数似然：
\begin{equation}
  \mathcal{L}_{\mathrm{GTS}}=-\sum_{(P,T)}\sum_{k=1}^{m}\log p\bigl(y_k\mid\bm{q}_k,\bm{c}_k\bigr),
  \label{eq:nlp-execution-gts-loss}
\end{equation}
其中$m$是树$T$的节点数，$y_k$是前序第$k$个节点的标注词元，$\bm{q}_k$与$\bm{c}_k$是该节点的目标向量和上下文向量。教师强制在这里有一个树上的含义：右子目标用到的$\bm{t}_{\mathrm{l}}$由标注的左子树合成，而不是由模型自己的预测合成。题中同一个数出现两次时，两个位置算作不同的候选词元，原文取概率较高的那个作为训练目标。推断时用一个栈保存尚待实现的目标，算法~\ref{alg:nlp-execution-gts}给出贪心版本；原文实际使用宽度为5的束搜索，束中的每个候选各自维护一个栈。

\begin{algorithm}[GTS目标驱动的表达式树解码（贪心版本）]
  \label{alg:nlp-execution-gts}
  \begin{algorithmic}[1]
    \Require 题目$P$的编码器状态$\bm{h}_{1:n}$，根目标$\bm{q}_0$，解码词表$V^{\mathrm{dec}}$
    \Ensure 表达式树$T$
    \State 把根节点及其目标$\bm{q}_0$压入目标栈$S$
    \While{$S$非空}
      \State 弹出栈顶节点$v$及其目标$\bm{q}$；用$\bm{q}$对$\bm{h}_{1:n}$做注意力，得到$\bm{c}$
      \State 按式\eqref{eq:nlp-execution-gts-score}选出$\hat y=\argmax_{y}p(y\mid\bm{q},\bm{c})$，令$v$的词元为$\hat y$
      \If{$\hat y\in V_{\mathrm{op}}$}
        \State 按式\eqref{eq:nlp-execution-gts-left}计算$\bm{q}_{\mathrm{l}}$；为$v$创建左、右子节点，右子节点标记“待左子树完成”
        \State 把右子节点、左子节点（含$\bm{q}_{\mathrm{l}}$）依次压入$S$ \Comment{左子节点先出栈}
      \Else
        \State 令$\bm{t}_v\gets\bm{e}(\hat y)$；沿祖先向上，凡左右子树均已完成的运算符节点，按式\eqref{eq:nlp-execution-gts-comb}合成其$\bm{t}$
        \State 若栈顶是左兄弟刚完成的右子节点，按式\eqref{eq:nlp-execution-gts-right}用该左兄弟的$\bm{t}_{\mathrm{l}}$计算其目标$\bm{q}_{\mathrm{r}}$
      \EndIf
    \EndWhile
    \State \Return $T$（按前序记录的词元，可直接求值）
  \end{algorithmic}
\end{algorithm}

图~\ref{fig:nlp-execution-gts}画出金额题的生成过程。第1步根目标预测出“$+$”，分解出左子目标；第2、3步左子目标继续预测出两个“$\times$”；第4、5步依次预测叶$n_1$与$n_0$，于是第3步的子树“$n_1\times n_0$”完成，它的嵌入$\bm{t}_{\mathrm{l}}$被用来生成第2步节点的右子目标；第6步该右子目标预测出$n_4$，第2步的子树“书的折后金额”完成，其嵌入再用来生成根的右子目标；第7至9步生成“$n_3\times n_2$”。九步之后栈空，前序序列为$+\ \times\ \times\ n_1\ n_0\ n_4\ \times\ n_3\ n_2$，代回数值即得式\eqref{eq:nlp-execution-price}的80元。整个过程中每个运算符恰好得到两个子节点，生成的序列总能还原成一棵合法的树，不存在括号失配。

\begin{figure}[!t]
  \centering
  % 第九卷 14.1.1.1：GTS目标驱动树解码在金额题上的生成过程（正文栏宽，约105 mm）。
% 表达式树 Add(Mul(Mul(25,3),0.8), Mul(4,5))，数量按原文顺序记为 n0=3, n1=25, n2=5, n3=4, n4=0.8。
% 节点左上角数字为前序生成的步次；蓝色实线：父目标向下分解出左、右子目标；
% 红色虚线：左子树完成后自底向上合成的子树嵌入 t，送去生成右兄弟的目标。
% FigureModel：运算符节点；FigureInput：数量叶节点。
\begin{tikzpicture}[
  x=1mm,y=1mm,
  figure picture,
  op/.style={figure node,draw=FigureModel,fill=FigureModelFill,minimum width=9mm,
    minimum height=6.5mm,inner sep=1pt},
  leaf/.style={figure node,draw=FigureInput,fill=FigureInputFill,minimum width=11mm,
    minimum height=6.5mm,inner sep=1pt},
  goal/.style={figure flow,draw=FigureInput},
  emb/.style={figure flow,draw=FigureOutput,dash pattern=on 2pt off 1.5pt},
  stepno/.style={figure note,inner sep=0.5pt}
]
  \node[op] (r) at (56,42) {$+$};
  \node[op] (a) at (28,29) {$\times$};
  \node[op] (b) at (84,29) {$\times$};
  \node[op] (c) at (14,16) {$\times$};
  \node[leaf] (n4) at (42,16) {$n_4{=}0.8$};
  \node[leaf] (n1) at (5,3) {$n_1{=}25$};
  \node[leaf] (n0) at (23,3) {$n_0{=}3$};
  \node[leaf] (n3) at (75,16) {$n_3{=}4$};
  \node[leaf] (n2) at (93,16) {$n_2{=}5$};

  \foreach \n/\k in {r/1,a/2,c/3,n1/4,n0/5,n4/6,b/7,n3/8,n2/9}{
    \node[stepno,anchor=south east] at (\n.north west) {\k};
  }

  \draw[goal] (r) -- node[figure edge label,above left=-1pt] {$\bm{q}_{\mathrm{l}}$} (a);
  \draw[goal] (r) -- node[figure edge label,above right=-1pt] {$\bm{q}_{\mathrm{r}}$} (b);
  \draw[goal] (a) -- (c);
  \draw[goal] (a) -- (n4);
  \draw[goal] (c) -- (n1);
  \draw[goal] (c) -- (n0);
  \draw[goal] (b) -- (n3);
  \draw[goal] (b) -- (n2);

  \draw[emb] (c.east) to[bend right=25] node[figure edge label,below=1pt] {$\bm{t}_{\mathrm{l}}$} (n4.west);
  \draw[emb] (a.east) to[bend right=12] node[figure edge label,below=1pt] {$\bm{t}_{\mathrm{l}}$} (b.west);

  \node[figure note,anchor=west,align=left] at (60,4)
    {前序：$+\ \times\ \times\ n_1\ n_0\ n_4\ \times\ n_3\ n_2$};
\end{tikzpicture}

  \caption{GTS在金额题上的目标驱动解码。节点旁的数字是前序生成的步次；蓝色实线表示父目标向下分解为左、右子目标$\bm{q}_{\mathrm{l}},\bm{q}_{\mathrm{r}}$；红色虚线表示左子树完成后合成的子树嵌入$\bm{t}_{\mathrm{l}}$，它与父信息一起生成右兄弟的目标。数量按原文顺序编号，$n_4=0.8$由“八折”规范化得到。}
  \label{fig:nlp-execution-gts}
\end{figure}

与序列到序列求解相比，GTS改变的是解码器的结构，编码器与训练数据都相同。它用树结构保证输出合法，用目标分解把“当前在为哪个子量服务”显式地传给每个节点，用子树嵌入让右侧的生成知道左侧已经算了什么。原文在Math23K数据集上做了一项对照：去掉子树嵌入、让右子目标与左子目标同样只依赖父信息时，五折交叉验证的答案正确率从74.3\%降到70.0\%，说明自底向上的信息对右子目标的生成确有作用。从更一般的角度看，GTS是第\ref{sec:nlp-inference-grammar}小节文法约束解码的一个特例：表达式的文法只有“运算符带两个子式”和“数量是叶”两条产生式，树解码器把这条文法写进了网络结构。它的局限也来自这一设定：只能生成单个表达式，不能表达方程组、多步赋值或需要循环的计算；数量必须先被正确识别为占位词元；答案正确率只核对数值，不同的树可能碰巧算出同一个数。后两类问题要由下文的程序执行与结果核验处理。

\subsubsection{有限域约束与回溯求解}
\label{sec:nlp-execution-csp}

排班题不能总用一个算式直接算出答案。\term{约束满足问题}（constraint satisfaction problem, CSP）为每个变量指定有限的候选域，再要求这些变量的联合取值满足全部条件。构造一个两时段值班例子：甲、乙、丙各选时段1或2；甲乙不能同时值班，乙必须早于丙，丙只能选时段2。令变量$a,b,c$分别表示三人的时段，便得到
\[
 D_a=D_b=\{1,2\},\quad D_c=\{2\},\qquad a\ne b,\quad b<c.
\]
本题求所有可行赋值，没有额外的最小成本目标。语言到约束的映射要逐词核对：若用布尔变量$z_i$表示人员是否被选择，“至少一人”对应$\sum_i z_i\geq1$，“恰好一人”对应$\sum_i z_i=1$，“甲乙不能同时被选”对应$z_{\mathrm{a}}+z_{\mathrm{b}}\leq1$，三者限定的是不同的集合。

\term{前向检查}（forward checking）是在每次赋值以后，删除尚未赋值变量中与当前赋值冲突的取值。固定变量顺序$a,b,c$，域内按升序尝试，并为每个分支保存一份域快照。令$a=1$后，$D_b$只剩$\{2\}$；再令$b=2$，约束$b<c$会把$D_c=\{2\}$删空。搜索回到$a$处，恢复原来的候选域，改试$a=2$，此时$D_b=\{1\}$，随后$b=1,c=2$满足全部条件。包括失败尝试在内，共进行了$a=1,b=2,a=2,b=1,c=2$五次赋值。更强的一致性传播可能更早发现$b=2$不可行，五次并不是这个问题的固有成本。

三种结论需要三种证据。找到$(2,1,2)$只说明问题可行；继续枚举完所有分支才能确认它是唯一解；去掉$a\ne b$后$(1,1,2)$也会成为解，而增加$a<b$则无解，此时要用穷尽的搜索树说明每个分支都已失败。若只找到一个解便停止，只能报告“找到可行排班”。

24点可以用类似的状态搜索，只是状态不再是人员取值。给定数字多重集$\{3,3,8,8\}$，每个数都须使用且只能使用一次，允许二元加、减、乘、除和括号。状态保存“剩余值、产生该值的表达式、已消耗的原始位置集合”。每次选两个不重叠的表达式合为一个值；减法与除法要保留两个顺序，零不能作除数。用有理数保存$8/3$，可以避免浮点误差把最终的$8/(3-8/3)=24$判成近似失败。

\subsubsection{模理论约束与Z3接口}
\label{sec:nlp-execution-z3}

约束涉及整数或实数算术与布尔组合时，手写搜索很快变得繁琐。\term{可满足性模理论}（satisfiability modulo theories, SMT）求解器判断是否存在某种解释，同时满足逻辑连接词和所用数域的规则。Z3是这种求解器的一个实现。下面用官方指南的SMT-LIB命令表达同一个排班题：声明类型、加入断言、检查可满足性，再读取模型。

\begin{codeblock}[label={code:nlp-execution-z3}]{排班约束的SMT-LIB表示}
(declare-const a Int)
(declare-const b Int)
(declare-const c Int)
(assert (and (<= 1 a) (<= a 2)))
(assert (and (<= 1 b) (<= b 2)))
(assert (= c 2))
(assert (distinct a b))
(assert (< b c))
(check-sat)
(get-model)
\end{codeblock}

这里\code{Int}限定整数，排除了1.5这样的中间时段。按前面的枚举，公式应为可满足，模型只能给出$a=2,b=1,c=2$。调用器只有在收到\code{sat}后才读取模型；\code{unsat}表示提交的公式没有模型，\code{unknown}表示这次检查未能作出决定，不能归为无解。要确认唯一性，可以加入$\lnot(a=2\land b=1\land c=2)$再检查一次，得到\code{unsat}即说明没有别的解；\code{push}/\code{pop}可以临时加入这类条件、检查后恢复原断言集。

求解器只对收到的公式负责。代码~\ref{code:nlp-execution-z3}若遗漏\code{distinct}，求解器会允许甲乙同时值班，它没有办法替作者补上这条限制。因此即使返回\code{sat}，也应把模型重新代回原题的每一条语言条件；实际使用时还应固定Z3版本并保存完整输出。与有限域搜索相比，SMT把搜索与传播交给经过长期优化的求解器，能处理大得多的整数与实数约束；代价是映射到SMT-LIB的这一步仍由模型或作者完成，而这一步正是语言任务中最容易出错的地方。

\subsection{程序辅助与工具集成推理}
\label{sec:nlp-execution-program-aided}

表达式树和约束都要求问题能写成某种固定的形式。更多题目需要多步计算、中间变量、循环或调用数学库，这时可以让语言模型直接写程序。纯文本推理在长计算中容易出算术错误，例如把$75\times0.8$算成65；程序辅助把算术交给解释器，解决的是计算的可靠性。但一次写完的程序无法根据中间结果调整：方程有两个根时该取哪个、数值解是否符合题意，都要看到执行结果后才能判断。工具集成推理让文本推理与代码执行交替进行，模型可以根据执行输出继续推理或改写代码。逐步写出推理文字的思维链、对多条路径投票的自一致性以及程序思维提示，已在第\ref{chap:nlp-reasoning}章讨论（第\ref{sec:nlp-execution-cot}、\ref{sec:nlp-execution-self-consistency}与\ref{sec:nlp-reasoning-program-aided}小节），本小节只展开程序执行怎样接入求解。

\subsubsection{程序辅助求解：PAL}
\label{sec:nlp-execution-pal}

第\ref{sec:nlp-reasoning-program-aided}小节介绍了程序辅助语言模型（PAL）的思想：模型把题目的中间关系写成程序，数值运算交给解释器\scite{nlp-execution-gao2023pal}。这里补充它接入求解时的实现。PAL只需要少样本提示：每个示例由题目和对应的Python程序组成，变量名带有语义（如\code{book\_price}），推理说明写在注释中。模型为新题生成程序以后，系统执行一次，从约定的变量或标准输出提取答案。模型负责“把语言条件翻译成程序”，解释器负责“按程序计算”，两者的分工与表达式树求解相同，只是中间表示从固定文法的树换成了通用程序。

代码~\ref{code:nlp-execution-pal}是按金额题构造的程序候选，已在本地Python 3解释器中实际执行，没有调用语言模型。它把折扣显式绑定到书的金额上，并用十进制定点数避免二进制浮点的显示误差。

\begin{codeblock}[language=Python,label={code:nlp-execution-pal}]{金额题的程序候选}
from decimal import Decimal

book_price = Decimal("25")
book_count = 3
book_pay_ratio = Decimal("0.8")
pen_price = Decimal("4")
pen_count = 5
book_payment = book_price * book_count * book_pay_ratio
pen_payment = pen_price * pen_count
answer = book_payment + pen_payment
print(f"{answer:.2f}")
\end{codeblock}

实际标准输出为\code{80.00}。结果提取器检查程序正常结束、输出只有一个合法金额，再附上原题单位形成“应付80元”。执行异常、数量绑定错误与折扣作用域错误是三类不同的失败，后两类程序都正常结束；区分它们需要保存程序、解释器版本、退出码与输出，见第\ref{sec:nlp-reasoning-program-aided}小节。生成的代码还应在限制权限、时间与资源的隔离环境中执行。

PAL消除了算术错误，却保留了两个局限。其一，程序在执行前一次写完，模型看不到任何中间结果；若题目需要先求出一个量、再根据它的取值决定下一步，单次程序只能把所有分支都预先写进去。其二，PAL把全部推理压进代码与注释，题意分析、分类讨论这类更适合用文字表达的步骤也被迫写成程序。

\subsubsection{代码与文本交替的工具集成推理}
\label{sec:nlp-execution-tora}

\term{工具集成推理}（tool-integrated reasoning）让模型在自然语言推理与程序执行之间来回切换。ToRA是一个代表\scite{nlp-execution-gou2024tora}。它的轨迹由三类片段交替组成：
\[
  \tau=r_1\,a_1\,o_1\,r_2\,a_2\,o_2\cdots r_{n-1}\,a_{n-1}\,o_{n-1}\,r_n,
\]
其中$r_i$是自然语言推理，$a_i$是一段程序，$o_i$是工具执行$a_i$的输出，最后的$r_n$在约定位置（如\code{\textbackslash boxed\{\}}）给出答案。模型生成$r_i$与$a_i$，系统执行$a_i$得到$o_i$并写回上下文，模型再以全部历史为条件生成下一段推理。生成过程在两类标记处切换：程序块写完、模型输出代表“请给出执行结果”的标记（原文为Markdown风格的\code{output}代码块起始标记）时，系统暂停生成、执行程序、把输出接在后面；推理中出现答案标记时，整条轨迹结束。算法~\ref{alg:nlp-execution-tora}给出这一循环。

\begin{algorithm}[工具集成推理的交替生成]
  \label{alg:nlp-execution-tora}
  \begin{algorithmic}[1]
    \Require 题目$q$，模型$G$，提示$p$（微调后可为空），执行器$E$，停止判据$\operatorname{Stop}(\cdot)$，最大轮数$n$
    \Ensure 轨迹$\tau$
    \State $\tau_0\gets$空串
    \For{$i=1,\ldots,n$}
      \State $r_i\sim P_G(\cdot\mid p\oplus q\oplus\tau_{i-1})$ \Comment{自然语言推理}
      \If{$\operatorname{Stop}(r_i)$} \Comment{$r_i$中已给出答案}
        \State \Return $\tau_{i-1}\oplus r_i$
      \EndIf
      \State $a_i\sim P_G(\cdot\mid p\oplus q\oplus\tau_{i-1}\oplus r_i)$ \Comment{生成到执行标记为止}
      \State $o_i\gets E(a_i)$ \Comment{执行程序，异常信息也作为输出}
      \State $\tau_i\gets\tau_{i-1}\oplus r_i\oplus a_i\oplus o_i$
    \EndFor
    \State \Return $\tau_n$ \Comment{达到轮数上限，可能没有答案}
  \end{algorithmic}
\end{algorithm}

用一道教学题说明交替的作用：“一个正数的平方比它的3倍大10，求这个数。”模型先写$r_1$：“设该数为$x$，则$x^2-3x-10=0$，用求根公式计算。”再写程序$a_1$，以\code{a, b, c = 1, -3, -10}求判别式$49$、整数平方根7，输出两根；该程序在本地Python 3中实际执行，输出$o_1$为\code{[-2, 5]}。接着模型读到两个根，写$r_2$：“题目要求正数，舍去$-2$，答案为$\boxed{5}$。”舍根这一步依赖执行结果，PAL只能把“取正根”预先写进程序，而且要事先知道会出现负根；工具集成推理则在看到$o_1$以后用文字处理它。程序报错时，错误信息同样作为$o_i$写回，模型可以在下一轮改写程序。

训练分两步。第一步是模仿学习：用GPT-4在GSM8K与MATH的训练题上按算法~\ref{alg:nlp-execution-tora}生成轨迹（先贪心解码，失败的题再做核采样），只保留答案正确且没有工具错误的轨迹，构成约1.6万条的ToRA-Corpus。在这些轨迹上微调开源模型，目标函数为
\begin{equation}
  \mathcal{L}_{\mathrm{IL}}=-\sum_{(q,\tau)}\sum_{i=0}^{n-1}\log P_{M}\bigl(r_{i+1}\,a_{i+1}\mid q,\,r_1a_1o_1\cdots r_io_i\bigr),
  \label{eq:nlp-execution-tora-loss}
\end{equation}
即只在模型自己要生成的推理与程序上计算损失；执行输出$o_i$只作为条件出现，不计入损失，因为它由解释器给出，模型不该学着“编造”执行结果。这与第\ref{sec:nlp-sft-masked-loss}小节对提示部分做损失掩码是同一种做法。

第二步是\term{输出空间塑形}（output space shaping）。ToRA-Corpus中每道题大多只有一条有效轨迹，模仿学习得到的模型输出方式单一，测试时难以探索别的合理解法。塑形分三步进行。先用模仿学习后的模型对每道训练题核采样64条轨迹，保留答案正确、无工具错误的有效轨迹。再利用错误轨迹：错误的解答往往在中途才出错，前半段仍然合理。把错误轨迹按行写成$\tilde\tau=l_1\cdots l_m$，枚举前缀$\tilde\tau[:j]=l_1\cdots l_j$，让一个更强的教师模型（原文为在ToRA-Corpus上训练的CodeLLaMA-34B）从该前缀出发贪心续写，答案正确的续写便是一条“被纠正”的轨迹。最后每题至多随机取4条采样或纠正得到的轨迹，与ToRA-Corpus合并，重新训练模型。纠正轨迹教给模型的是“从一个看似合理但会出错的位置怎样走对”，这正是只看示范时学不到的。

ToRA与前文方法的关系可以这样看。PAL是ToRA在一轮、且$r_1$为空时的特例；ToRA的“采样—按答案筛选—再训练”与第\ref{sec:nlp-reasoning-rft}小节的拒绝采样微调同属一类，区别在于轨迹中含有真实的执行输出，并多了从错误前缀出发的纠正；它的推理—行动—观察交替又与第\ref{sec:nlp-execution-environment}节的ReAct同构，只是工具固定为解释器，任务限于数学。ToRA的优点是计算可靠、能依据中间结果调整，并且训练后不再依赖长的少样本提示。代价是每一轮都要启动执行，推断延迟随轮数增加，原文把工具执行次数上限设为3；训练数据依赖强教师模型与带标准答案的题目；筛选只看最终答案，答案对而程序逻辑错的轨迹仍会进入训练集，这要由下一小节的核验处理。

\subsection{结果验证与解释}
\label{sec:nlp-execution-verification}

候选形成后，需要分别检查最终答案、中间过程和解释三者是否成立并互相对应。一条解答可能碰巧给出80元，却包含错误等式；也可能程序算对了，随后的文字却把折扣对象说错。检查手段可按可靠性与覆盖面排列：独立规则回代最可靠，但只覆盖能写出规则的题；学习得到的奖励模型能给任意候选打分，用于在多个候选中选择或作为训练信号，但它的判断本身是预测。

\subsubsection{计算回代与规则核验}
\label{sec:nlp-execution-rule-check}

规则核验器读取原题的变量表、允许运算与目标，独立重算候选。金额题可用以分为单位的整数关系交叉检查：书支付$2500\times3\times4/5=6000$分，笔支付$400\times5=2000$分，总计8000分。核验必须从原题重新绑定折扣范围，并比较书、笔两项明细；若与PAL程序共用同一段错误公式，只会重复同一个错误。

24点的检查分成语法、数值与资源使用三项。解析器只允许整数叶与四则二元运算，从树中收集叶值及出现次数，检查是否恰为$\{3,3,8,8\}$，再以有理数递归求值并拒绝除零。表达式$8\times3+3-3$的值确实是24，却用了三个3、一个8，应被拒绝；直接提交常数24也没有消耗题给的数字。搜索器保存的位置集合便于构造，独立核验器则从最终表达式重建证据，不信任搜索过程中的状态记录。

规则核验要预先约定数域、单位、舍入与容差，以及题目是否允许多个答案；否则“数值不同”可能只是表示不同，“数值相同”也可能掩盖对象不同。

\subsubsection{结果与过程奖励模型的应用}
\label{sec:nlp-execution-reward-models}

许多自然语言解答不能完全转成规则检查，这时可以用学习的评分器帮助选择。结果奖励模型（ORM）读取题目与完整解答，预测该解答最终正确的概率；过程奖励模型（PRM）在每个推理步骤末尾预测这一步是否正确。二者的结构、训练目标（式\eqref{eq:nlp-reasoning-orm}与式\eqref{eq:nlp-reasoning-prm}）、人工步骤标注与蒙特卡洛自动标注以及步骤分数的聚合规则，已在第\ref{sec:nlp-reasoning-orm}至\ref{sec:nlp-reasoning-prm-aggregation}小节讲过。这里只说明在数学求解任务中怎样构造标签、怎样使用它们。

数学题的优势是标签可以自动得到。ORM的标签来自最终答案：从解答中按约定位置提取答案，与标准答案做等价判断。等价判断本身就是一个规则核验器，其规范化要求见第\ref{sec:nlp-reasoning-rewards}小节。PRM的步骤标签可以用第\ref{sec:nlp-reasoning-prm-mc}小节的蒙特卡洛续写估计：从某一步的前缀出发续写多次，以到达正确答案的比例作为这一步的软标签。训练、调参与测试应按题目划分，并记录候选来自哪个生成器，否则验证器可能只是记住了训练题或熟悉的错误风格。

在多候选选择中，两种模型解决的问题不同。设同一道金额题采样了六条解答，五条把笔也打了折、给出76元，一条给出80元。多数投票会选76元；若ORM学会了识别“折扣作用域错误”这类模式，它可以把正确的那一条排在首位。评价时应分开统计三个量：候选集合中是否有正确解，正确候选能否排在错误候选之前，以及最终选出的那一条是否正确。前一个量不足要改进生成，后两个量不足要改进评分器。

ORM的标签只看答案，会把“步骤错误但末尾碰巧写对”的候选标成正例。构造三步候选：“书原价75元；$75\times0.8=65$元；加20元，合计80元”。它的答案正确，第二步乘法和第三步加法却都有错误，ORM的标签会把它当成正例。PRM在步骤上评分，若三步得分为$0.99,0.10,0.98$，按乘积聚合为$0.09702$；另一条正确的三步解答每步得0.9，乘积为$0.729$，排序便偏向后者（数值为手设的教学示例）。第一个低于阈值的步骤给出了复核位置，但确认$75\times0.8\ne65$仍应交给算术检查。PRM还可以放入搜索，在每一步只保留评分高的前缀，见第\ref{sec:nlp-reasoning-prm-beam}小节；这时每次前缀评分都要计入调用成本。两类模型都只是预测，它们在策略分布变化后可能失准，可以规则核验的部分仍应交给规则。

\subsection{形式化定理证明}
\label{sec:nlp-execution-theorem-proving}

证明无法用数值回代检查，自然语言证明的每一步又可能含有不严密的跳跃。\term{证明助手}（proof assistant）把命题和证明都写成形式语言中的对象，由一个很小的可信内核逐步检查。只要内核本身正确，通过检查的证明在所用的公理系统内就是正确的，验证不再依赖评分器的判断。这把困难转移到了两处：把自然语言命题准确地写成形式命题，以及在巨大的证明空间中找到一个能通过检查的证明。本小节先说明Lean怎样检查证明，再讨论用语言模型自动形式化命题，最后讨论用语言模型生成证明步骤并配合搜索。

\subsubsection{Lean证明对象检查}
\label{sec:nlp-execution-lean}

Lean 4把命题表示为类型，把证明表示为该类型的一个项：若能构造出类型为$P$的项，$P$就得到证明。检查证明于是变成检查项的类型，这由内核完成。下例只使用命题、函数应用和合取类型，不依赖额外的数学库。

\begin{codeblock}[label={code:nlp-execution-lean}]{从已知前提构造证明项}
theorem eligibility
    (registered timely allowed : Prop)
    (rule : And registered timely -> allowed)
    (facts : And registered timely) : allowed :=
  rule facts

#print axioms eligibility
\end{codeblock}

\code{registered}表示已提交信息，\code{timely}表示在截止时刻前提交，\code{rule}是明确给出的前提“同时满足两条件则允许”。\code{rule}的类型是从合取到\code{allowed}的函数，把\code{facts}交给它，得到的项类型恰为\code{allowed}，于是证明完成。验收时要固定Lean版本与依赖库，排除\code{sorry}、\code{admit}等未完成占位，并审查\code{\#print axioms}列出的公理依赖；文件没有语法错误还不足以说明证明完成。

直接写证明项对复杂命题并不现实，实际证明多用\term{证明策略}（tactic）书写。证明策略是一条改变证明目标的指令：系统维护当前的\term{证明状态}（proof state），即若干待证目标，每个目标由已知假设和一个待证结论组成；执行一条策略后，当前目标被关闭或被替换为新的子目标，所有目标都关闭时证明完成。同一定理用策略写成：
\begin{codeblock}[label={code:nlp-execution-lean-tactic}]{用证明策略书写同一证明}
theorem eligibility'
    (registered timely allowed : Prop)
    (rule : And registered timely -> allowed)
    (facts : And registered timely) : allowed := by
  apply rule
  exact facts
\end{codeblock}
初始状态的目标是\code{allowed}；\code{apply rule}用\code{rule}的结论匹配目标，把目标换成它的前提\code{registered}$\land$\code{timely}；\code{exact facts}用已有假设关闭这个目标。策略程序最终仍被展开为证明项交给内核检查，因此策略本身写错只会导致检查失败，不会让错误的证明通过。两段代码只用到Lean 4核心库的定义与策略，正文未在本地编译，使用时应记录\code{lean --version}并保存检查器输出。证明策略的这种“状态—指令—新状态”结构，正是下文让语言模型逐步生成证明的基础。

形式证明只对提交的命题负责。若把\code{timely}错误地定义为“4月16日当天”，证明仍会通过，却只证明了这个错误的定义，原通知中“4月16日18:00前”的要求没有被检查到。形式化映射与证明检查因此是两项独立的职责：前者对齐语言要求，后者确认在给定定义和前提下结论成立。

\subsubsection{自然语言命题的自动形式化}
\label{sec:nlp-execution-autoformalization}

人工把数学题写成形式命题需要熟悉证明助手与数学库的专家，形式命题的数量因此远少于自然语言题目，这直接限制了证明模型的训练数据。\term{自动形式化}（autoformalization）用语言模型把自然语言命题翻译为形式命题。Wu等较早系统地研究了这件事：用少样本提示让代码语言模型把竞赛数学题译为Isabelle/HOL命题，人工检查发现约四分之一（25.3\%）的译文完全正确；把译出的命题用于专家迭代训练证明器，也提高了它在miniF2F基准上的证明率\scite{nlp-execution-wu2022autoformalization}。此后的工作多在Lean 4上以微调的翻译模型实现，形式与机器翻译相同：输入自然语言题目，输出一条以\code{sorry}占位证明的形式命题。

自动形式化的难点是语义对齐。用一道教学题说明：“若实数$x$满足$x^2=4$且$x>0$，则$x=2$。”忠实的译文是\code{theorem t (x : Real) (h0 : x \textasciicircum{} 2 = 4) (h1 : 0 < x) : x = 2}（需要Mathlib中的实数）。常见的错误有三类。漏掉假设\code{h1}，命题变为假，证明搜索会失败，但失败本身不说明原因。把\code{Real}换成自然数\code{Nat}，命题仍然为真且容易证明，可它说的已是另一件事：在自然数中负根根本不存在，“$x>0$”这个条件被无声地吸收了；Lean中自然数的减法截断为零、除法向下取整，也常使译文与原意不符。第三类是假设互相矛盾，例如把条件误译为\code{h0 : x \textasciicircum{} 2 = 4}与\code{h1 : x < -3}，这时任何结论都能证明，证明成功反而毫无意义。

针对这三类错误，筛选分几层进行。编译检查确认译文在固定版本的Lean与Mathlib中能通过类型检查，它排除语法与类型错误，却发现不了上述任何一类语义偏差。回译检查把形式命题再译回自然语言，与原题比较是否表达同一个命题；Lean Workbook用编译器与自然语言推断模型判断回译与原题是否一致，并把不一致的样例交给专家修正、加入训练集，迭代改进翻译模型\scite{nlp-execution-ying2024leanworkbook}。DeepSeek-Prover在模型评分之外加入\term{假设拒绝}（hypothesis rejection）：尝试从命题的假设证明\code{False}，若能证出，说明假设矛盾，该命题被丢弃；它还对命题与其否定同时搜索证明，若证出否定，说明译文为假\scite{nlp-execution-xin2024deepseekprover}。这些检查从不同方向补充了编译检查，但都不完备：回译模型可能把两处偏差同时“修正”回原意，自然数替换实数这样的错误也未必能被自然语言推断发现。形式定义本身错了，内核无从得知，最终仍要对代表性样本做人工审查。

\subsubsection{证明策略生成与搜索}
\label{sec:nlp-execution-proof-search}

有了形式命题，下一步是找证明。证明策略的结构把证明变成一个搜索问题：状态是证明状态，动作是一条策略，Lean执行策略给出新状态或报错，所有目标关闭即到达终点。可选的策略及其参数组合极多，盲目枚举不可行；语言模型的作用是为当前状态提出少量有希望的策略。

\paragraph{策略语言模型}
GPT-f最早用语言模型生成证明步骤：把证明状态序列化为文本，训练模型在“目标$\to$下一步”的格式上做条件语言建模\scite{nlp-execution-polu2020gptf}。原文的训练样本形如\code{GOAL <目标> PROOFSTEP <步骤>}，推断时给出\code{GOAL <目标> PROOFSTEP}，让模型续写步骤。训练数据从已有的形式数学库中抽取：把库中每个证明逐步执行，记录每条策略执行前的状态，得到大量“状态—策略”对，目标函数是策略词元的负对数似然，只在策略部分计算损失。GPT-f针对Metamath，后续工作把同样的格式用于Lean。

\paragraph{最佳优先搜索}
单次采样的一条策略序列很可能中途失败，需要在多个分支间分配尝试。\term{最佳优先搜索}（best-first search）维护一个按优先级排序的待扩展状态队列，每次取出优先级最高的状态扩展。GPT-f使用累计对数概率作为优先级：
\begin{equation}
  \operatorname{prio}(s)=\sum_{j=1}^{J}\log p_{\bm\theta}(a_j\mid s_{j-1}),
  \label{eq:nlp-execution-proof-priority}
\end{equation}
其中$s_0\xrightarrow{a_1}s_1\cdots\xrightarrow{a_J}s_J=s$是从初始状态到$s$的策略路径，$p_{\bm\theta}$是策略语言模型。优先级越高，表示模型越相信这条路径。由于每一步都加一个负数，较深的状态优先级往往较低，搜索因此带有广度优先的倾向。算法~\ref{alg:nlp-execution-best-first}以证明状态为节点给出完整过程。

\begin{algorithm}[基于策略语言模型的最佳优先证明搜索]
  \label{alg:nlp-execution-best-first}
  \begin{algorithmic}[1]
    \Require 形式命题对应的初始状态$s_0$，策略模型$p_{\bm\theta}$，每次扩展的采样数$e$，扩展次数上限$B$
    \Ensure 证明（策略序列）或“未找到”
    \State 优先队列$Q\gets\{(s_0,0)\}$；已访问集合$U\gets\{s_0\}$
    \For{$k=1,\ldots,B$}
      \If{$Q$为空} \Return 未找到 \EndIf
      \State 取出$Q$中优先级最高的$(s,\pi)$
      \State 由$p_{\bm\theta}(\cdot\mid s)$采样或束搜索得到$e$条不同的策略$a^{(1)},\ldots,a^{(e)}$
      \For{每条策略$a$}
        \State 在Lean中对$s$执行$a$；若报错或超时，丢弃
        \State 得到新状态$s'$；若$s'$没有剩余目标，\Return 从$s_0$到$s'$的策略序列
        \If{$s'\notin U$}
          \State 把$(s',\pi+\log p_{\bm\theta}(a\mid s))$加入$Q$，$s'$加入$U$，并记录父指针
        \EndIf
      \EndFor
    \EndFor
    \State \Return 未找到 \Comment{预算耗尽，不能据此断言命题不可证}
  \end{algorithmic}
\end{algorithm}

用代码~\ref{code:nlp-execution-lean-tactic}的定理走一遍（对数概率为教学假设，策略效果按Lean 4的语义给出）。第1次扩展取出初始状态$s_0$，目标为\code{allowed}，设模型提出三条策略：\code{apply rule}（$-0.4$）、\code{exact facts}（$-0.9$）、\code{constructor}（$-1.6$）。后两条在Lean中失败：\code{facts}的类型是合取而不是\code{allowed}，\code{constructor}要求目标是归纳类型。只有\code{apply rule}成功，得到状态$s_1$，目标为\code{registered}$\land$\code{timely}，优先级$-0.4$。第2次扩展取出$s_1$，设模型提出\code{exact facts}（$-0.1$）与\code{constructor}（$-0.7$）：前者关闭全部目标，搜索返回证明“\code{apply rule}；\code{exact facts}”，路径优先级为$-0.5$；后者会把状态拆成\code{registered}与\code{timely}两个目标，只是已不必再扩展。两次扩展、五次策略执行就找到了证明。真实问题中失败的分支远多于成功的分支，扩展上限$B$与采样数$e$决定了计算量。

\paragraph{价值估计}
累计对数概率只反映模型“多常这样写”，不反映“从这里还能不能证完”。GPT-f因此又训练了一个可证性函数：在搜索中访问过的目标，最终被证明的标为正，否则标为负，用同一个语言模型在\code{GOAL <目标> OUTCOME}之后预测\code{P}或\code{N}，以\code{P}的概率$f_{\mathrm{P}}(g)$作为目标$g$可证的估计。一条策略产生的子目标必须全部证明，所以目标的价值取同一策略产生的全部子目标（含$g$）的可证概率之积，$V(g)=\prod_{g'\in\operatorname{sib}(g)}f_{\mathrm{P}}(g')$；后续迭代用$V$代替累计对数概率作为优先级。这与第\ref{sec:nlp-reasoning-prm-mc}小节“从这一步出发还能做对的概率”是同一种价值，只是成功与否由证明助手判定。HTPS（HyperTree Proof Search）把这种“与”结构显式地写进搜索：它把搜索表示为超图，节点是单个目标，超边从一个目标指向策略产生的一组子目标，在其上做受AlphaZero启发、类似蒙特卡洛树搜索的选择与回传，并用评价模型估计目标可证的概率\scite{nlp-execution-lample2022htps}。语言状态上的蒙特卡洛树搜索见第\ref{sec:nlp-reasoning-mcts}小节；证明搜索的特殊之处在于“与”结构：父目标的价值取决于其全部子目标，而不是其中最好的一个。

\paragraph{专家迭代式训练}
搜索找到的证明可以反过来训练模型。Polu等在Lean上采用专家迭代：用当前模型对一批形式命题做证明搜索，把成功的证明拆成“状态—策略”对加入训练集，重新训练模型，再用新模型搜索，如此反复\scite{nlp-execution-polu2023curriculum}。搜索充当“专家”，模型充当“学徒”，这与第\ref{sec:nlp-reasoning-rft}小节推理轨迹上的专家迭代是同一框架，只是筛选器换成了证明检查，几乎没有假阳性。他们观察到，命题集合的难度分布足够分散时，迭代会自发形成课程：先证出的容易命题使模型变强，再证出稍难的命题。DeepSeek-Prover把这一循环与自动形式化结合：先把大量自然语言竞赛题自动形式化并筛选，再用模型生成证明，经Lean验证的证明加入训练集，更新后的模型再用于下一轮形式化与证明，迭代多轮\scite{nlp-execution-xin2024deepseekprover}。与逐条策略搜索不同，它让模型一次生成完整证明，再整体交给Lean检查；这样每次尝试只调用一次检查器，代价是失败时得不到中间状态的反馈，只能整体重采样。

三种做法的取舍在于“每次调用证明助手得到多少信息”。逐策略搜索每一步都得到新状态，能及时剪掉失败分支，但与证明助手的交互次数多；整体生成交互最少，依赖模型一次写对；价值估计与超图搜索把交互花在更有希望的分支上，代价是要额外训练评价模型。三者共享同一项保证：只要命题形式化正确，找到的证明就是正确的，这是数值回代与奖励模型都给不出的。搜索失败则只说明在给定预算内没有找到证明，既不说明命题为假，也不说明形式化有误。

\begin{table*}[t]
  \centering
  \small
  \caption{数学与逻辑问题求解的方法对照。“新代价”指方法在解决前一问题时引入的成本或残留风险；各方法可以组合使用，如工具集成推理产生的候选再由规则核验或奖励模型选择。}
  \label{tab:nlp-execution-solving-methods}
  \begin{tabularx}{\linewidth}{@{}p{30mm}p{34mm}Xp{40mm}@{}}
    \toprule
    方法 & 解决的问题 & 核心机制 & 新代价\\
    \midrule
    序列到序列求解 & 不用人工特征生成表达式 & 数量占位，编码器—解码器生成表达式序列 & 括号失配；不知道当前在为哪个子量服务\\
    GTS & 输出合法性与生成的结构信息 & 目标自顶向下分解，子树嵌入自底向上合成，前序生成 & 只能表达单个表达式；依赖数量识别\\
    CSP与SMT & 多变量条件的联合满足 & 回溯与传播；交给求解器判定sat/unsat & 约束映射仍由人或模型完成\\
    PAL & 长计算中的算术错误 & 生成程序，由解释器执行一次 & 看不到中间结果；需要隔离执行\\
    ToRA & 需要依据中间结果调整的推理 & 推理、程序、执行输出交替；模仿学习与输出空间塑形 & 多轮执行的延迟；依赖强教师与带答案的题\\
    规则回代 & 答案与资源使用的独立确认 & 从原题重建变量并重算 & 只覆盖能写出规则的题\\
    ORM与PRM & 不可规则化候选的选择与定位 & 按答案或步骤标签训练评分器 & 预测会失准；PRM标签成本高\\
    自动形式化 & 形式命题稀缺 & 翻译为Lean命题，编译、回译与假设拒绝筛选 & 语义偏差无法由内核发现\\
    证明搜索 & 在巨大策略空间中找证明 & 策略语言模型提议，最佳优先或树搜索，专家迭代 & 与证明助手的大量交互；失败不说明不可证\\
    \bottomrule
  \end{tabularx}
\end{table*}

\section{程序任务}
\label{sec:nlp-execution-programs}

第\ref{sec:nlp-execution-pal}小节的PAL把程序当作求解题目的中间工具，程序写错了，最终答案也会错，但程序本身不交付给用户。程序任务则直接交付软件产物或对其行为的说明：解释一个现有函数做了什么，按规格写出新函数，把代码迁到另一种语言或接口，或者修复仓库中的错误。验收依据也随之改变：数学题核对一个答案，程序要核对约定输入上的全部行为。

这类任务有四个挑战，本节各用一个小节处理：理解既有代码，函数名和注释只提供意图线索，真实行为要从跨文件的代码、语法结构和运行中恢复；生成正确代码，单次采样的正确率有限，输出须能被解析，功能错误只有执行才能暴露；在语言间转换并保持行为，平行代码稀少，两种语言的运算语义又有细微差别；在大型仓库中定位并修复问题，相关代码散落在许多文件中，修改须解决目标问题而不破坏其余行为。

方法谱系按“证据从哪里来、候选怎样形成、怎样筛选”展开。程序理解的证据从文本检索走向语法结构，再走向运行轨迹。代码生成由代码语言模型提供候选，再沿三条互补的路线提高交付质量：大规模采样后按测试和行为筛选，解码时用文法约束排除语法错误，生成后把执行结果写回提示让模型修订。代码转换从人工编写的语法树规则走向只用单语代码训练的回译式翻译。仓库级修复先有让模型自主浏览、编辑与测试的智能体，再有用固定流程替代自主决策的Agentless。表\ref{tab:nlp-execution-program-methods}按同一口径列出这些方法。

\begin{table*}[htbp]
  \centering
  \small
  \begin{tabularx}{\linewidth}{@{}l>{\raggedright\arraybackslash}X>{\raggedright\arraybackslash}X>{\raggedright\arraybackslash}X@{}}
    \toprule
    方法 & 解决的问题 & 核心机制 & 新代价或依赖 \\
    \midrule
    检索、AST与运行轨迹 & 从代码恢复当前行为 & 符号与语义检索取证据，语法树追踪控制流与数据流，执行记录实际返回 & 静态分析不覆盖动态调用，运行只覆盖所测输入 \\
    代码语言模型与FIM & 获得能写代码、能按光标两侧补全的模型 & 去重与许可过滤后的代码语料，中间填充重排，仓库级拼接 & 数据治理成本；FIM比例过高损害左到右续写 \\
    AlphaCode & 单次采样正确率低而提交次数有限 & 大量采样，示例测试过滤，在生成输入上按输出聚类后每类取一个 & 采样与执行的计算量；依赖生成输入区分行为 \\
    SynCode & 自由生成可能语法错误 & 增量解析得到接受序列，按终结符DFA状态查预计算的词表掩码 & 文法须为LR(1)/LALR(1)；只保证语法 \\
    Self-Debugging & 重采样丢弃了失败候选中的信息 & 测试结果、执行轨迹或代码解释写回提示，多轮修订 & 多轮调用；无测试时反馈来自模型自己 \\
    AST规则转换 & 精确地做接口迁移 & 匹配语法树节点并重写 & 规则须人工编写，覆盖有限 \\
    TransCoder & 缺少平行代码 & 跨语言掩码语言模型预训练、去噪自编码、回译 & 无行为保证，须用单元测试评价 \\
    SWE-agent & 模型直接用命令行浏览编辑仓库易出错 & 为模型设计的查看、搜索、带检查的编辑命令和简洁反馈 & 长轨迹的调用成本，错误会沿轨迹累积 \\
    Agentless & 自主智能体的决策复杂、成本高、难分析 & 固定的分层定位、局部修复与补丁验证流程 & 流程之外的情形无法自行调整 \\
    \bottomrule
  \end{tabularx}
  \caption{程序任务的方法谱系。“新代价或依赖”指相对于前一种做法新增的计算、数据或适用条件。}
  \label{tab:nlp-execution-program-methods}
\end{table*}

各方法共用一个前提：固定代码版本、运行环境、依赖和允许修改的范围，才能知道解释针对哪份程序、测试实际执行了什么。图\ref{fig:nlp-execution-program-loop}将开发反馈与独立验收分开。开发测试的结果可以进入候选修订；独立验收只能在候选冻结后评价，否则它也成为开发信息。

\begin{figure*}[htbp]
  \centering
  % Adapted from academic-drawing/templates/tikz_template.tex.
% Main ports: south -> north. Feedback remains left of the main column.
\begin{tikzpicture}[
  every node/.style={font=\figurefont\small,align=center,text=BookInk},
  block/.style={book node,font=\figurefont\small,text width=44mm,minimum height=11mm},
  flow/.style={book arrow,rounded corners=1.2mm},
  feedback/.style={flow,dashed,draw=BookAmber}
]
  \node[block,draw=FigureInput,fill=FigureInput!10] (spec)
    at (-45mm,0) {规格与允许变化\\输入、输出、保持条件};
  \node[block,draw=FigureInput,fill=FigureInput!10] (assets)
    at (45mm,0) {已有资产\\仓库版本、接口与依赖};
  \node[block,draw=FigureModel,fill=FigureModel!10] (candidate)
    at (0,-25mm) {生成候选或局部补丁\\保留版本与修改位置};
  \node[block] (test) at (0,-47mm)
    {在固定环境运行\\目标测试与相关回归};
  \node[block] (result) at (0,-69mm)
    {读取实际结果\\返回值、异常、失败断言};
  \node[block,text width=27mm,draw=BookAmber,fill=BookAmber!8]
    (revise) at (-61mm,-69mm) {定位并修订\\计入总预算};
  \node[block,draw=FigureOutput,fill=FigureOutput!10] (freeze)
    at (0,-91mm) {按预定规则选定\\冻结候选与环境};
  \node[block,text width=38mm] (stop) at (60mm,-69mm)
    {停止并报告\\预算耗尽或环境失败};
  \node[block,draw=FigureOutput,fill=FigureOutput!10] (heldout)
    at (0,-118mm) {独立验收\\留出测试与结果记录};
  \coordinate (inputs) at (0,-14mm);
  \draw[draw=BookMuted,line width=.7pt,rounded corners=1.2mm]
    (spec.south) |- (inputs);
  \draw[draw=BookMuted,line width=.7pt,rounded corners=1.2mm]
    (assets.south) |- (inputs);
  \draw[flow] (inputs) -- (candidate.north);
  \draw[flow] (candidate.south) -- (test.north);
  \draw[flow] (test.south) -- (result.north);
  \draw[feedback] (result.west) -- node[above] {未通过} (revise.east);
  \draw[feedback] (revise.north) |- (candidate.west);
  \draw[flow] (result.east) -- (stop.west);
  \draw[flow] (result.south) -- node[right] {通过} (freeze.north);
  \draw[flow] (freeze.south) -- (heldout.north);
  \draw[draw=BookMuted,line width=.8pt,dotted]
    (-78mm,-104mm) -- (80mm,-104mm);
  \node[anchor=east,align=right] at (79mm,-114mm)
    {不回流为\\开发反馈};
\end{tikzpicture}

  \caption{程序候选、实际执行与修订。实线给出从规格到验收的处理路径，虚线把开发失败反馈给修订；点线划开开发与独立评价。每个候选使用相同环境和测试协议，测试结论只覆盖约定的行为范围。}
  \label{fig:nlp-execution-program-loop}
\end{figure*}
\FloatBarrier

\subsection{程序理解与规格恢复}
\label{sec:nlp-execution-understanding}

“这个函数按通知发布时间排序吗？”不能只读函数名作答。程序理解先检索与问题相关的代码并连接跨文件依赖，再借语法结构和真实运行恢复当前行为。三类证据从文本相似到结构到运行逐步加强；由代码恢复的是当前行为的说明，它是否符合原始意图，还要对照需求、接口契约或有效测试等独立规格。

\subsubsection{符号检索与BM25、嵌入检索}
\label{sec:nlp-execution-code-retrieval}

对仓库提问时，检索单位宜是路径、限定符号名、函数签名、注释和代码片段组成的记录，而不是任意字符窗口。问题“通知排序为何忽略同一天内的先后”可以拆成符号查询\code{sort\_notices}以及“发布时间、日期解析、排序键”等词。三种检索互补：名称匹配取得直接定义，BM25按词项召回候选（打分公式见第\ref{sec:nlp-classification-bm25}小节），双塔嵌入检索补充措辞不同而语义相关的辅助函数（训练方法见第\ref{sec:nlp-classification-dpr}小节）。

设仓库同时有\code{notices.py:sort\_notices}和测试文件中的同名辅助函数，仅按名称会召回两者；检索器若经监督训练，这种同名干扰项正是有用的困难负例。读取生产入口的导入关系可以确认它调用了\code{time\_utils.py:parse\_local}，再扩展到该定义，就能看到时间字段是否在解析中被截断。交给模型的材料应带路径、行号和版本，并注明尚未读到的依赖。排名靠前的片段只回答“可能在哪里”，沿调用链取得的跨文件证据才开始回答“为什么会这样”。

\subsubsection{AST与控制流、数据流证据}
\label{sec:nlp-execution-ast}

取得代码后，需要追踪条件和数据怎样影响返回值。\term{抽象语法树}（abstract syntax tree, AST）把源码解析为函数、条件、调用、运算与返回等语法节点，Python标准库的\code{ast.parse}提供这一结构。对代码\ref{code:nlp-execution-sort-trace}，AST显示一个空输入分支、一个类型检查分支，以及由\code{set}传给\code{sorted}的嵌套调用。

\begin{codeblock}[language=Python,label={code:nlp-execution-sort-trace}]{需要核对去重行为的排序函数}
def unique_sorted(values):
    if not values:
        return []
    if any(type(x) is not int for x in values):
        raise TypeError("integers required")
    return sorted(set(values))
\end{codeblock}

在AST上连接条件的真假出口、返回与异常出口，得到控制流：空列表直接返回；非空输入检查全部元素类型，不合格时抛出异常。跟踪变量的定义与使用，得到数据流：\code{values}先进入集合构造，重复值在排序之前已被合并，注释若声称“保留重复项”就与实际计算不符。类型判断用\code{type(x) is not int}，所以布尔值也被拒绝，尽管Python中布尔类型是整数的子类。导入名被重新赋值或调用目标来自运行时注册表时，静态调用边只能标为“可能”，这类不确定性应写进解释的限制。

\subsubsection{运行轨迹与边界输入}
\label{sec:nlp-execution-traces}

把具体输入绑定到程序并执行，可以记录实际返回、异常与关键中间变量。代码\ref{code:nlp-execution-sort-trace}对空列表返回\code{[]}，对\code{[3,1,3]}经集合\code{\{1,3\}}返回\code{[1,3]}，对\code{[1,"2"]}在类型检查处抛出\code{TypeError}。三个输入分别走过空分支、重复项合并和异常分支，使“排序并去重、只接受整数”的说明落到具体代码位置。

轨迹须来自运行时采集（脚本捕获的返回值与异常、调试器中的变量快照），并保留环境版本；模型写出的“第三行得到两个元素”只能引导核查。三个样本没有覆盖生成器、自定义容器等输入，分支覆盖也不衡量输入空间是否穷尽，所以恢复的行为说明只适用于所测输入范围。

\subsection{代码生成}
\label{sec:nlp-execution-generation}

代码生成以规格为输入、以程序为输出。以通知排序为例，规格约定输入为记录列表，每条记录含唯一业务标识和带时区的发布时间；输出按时间升序，保留重复记录与相同时刻记录的原次序，不修改输入；缺失时区、无效日期和不合格记录应报错而非猜测。正常、边界与无效输入的测试预期都从这份规格推出，不从候选代码反推。

这里有三个挑战：单次采样的正确率有限，正确程序往往要在许多候选中才找得到；输出须符合目标语言的语法，否则无法运行；功能错误多数只有执行才能暴露。下面先说明候选来自怎样训练的代码语言模型，再依次讨论针对这三个挑战的采样筛选、文法约束与执行反馈修订。多候选方法常用pass@$k$评价（无偏估计见式\eqref{eq:nlp-eval-passk}），它衡量候选覆盖，不衡量从候选中选对的能力，后者正是筛选与修订要解决的问题。

\subsubsection{代码语言模型的训练数据与目标}
\label{sec:nlp-execution-code-lm}

通用语言模型的语料里代码占比有限，而写代码需要熟悉大量库接口、语法细节和跨文件引用。\term{代码语言模型}（code language model）用以代码为主的语料预训练或继续训练，结构仍是第\ref{chap:nlp-large-language-models}章的仅解码器Transformer，主要差别在数据和训练样本的构造。Codex在GPT-3基础上用公开代码仓库继续训练\cite{nlp-intro-chen2021}，StarCoder、DeepSeek-Coder等开放模型公开了更完整的数据流程\cite{nlp-execution-li2023starcoder,nlp-execution-guo2024deepseekcoder}。下面依次说明数据怎样处理、目标怎样改造、仓库结构怎样进入样本。

数据处理在第\ref{sec:nlp-llm-data}节的通用流程上增加了代码特有的步骤。第一步是许可证筛选：The Stack只收录宽松许可证（如MIT、Apache-2.0）下的公开代码，并提供让作者查询和申请删除自己代码的退出机制\scite{nlp-execution-kocetkov2022stack}。第二步是去重：开源代码中复制、派生和自动生成的文件极多，按第\ref{sec:nlp-llm-data-minhash}小节的MinHash做近重复去除，The Stack的实验中近重复去除在全部设置下都提高了下游效果。第三步是质量过滤：删去行过长、字母字符比例过低或明显自动生成的文件，并移除密钥、邮箱等个人信息。第四步是去污染：评测题目和参考解常被上传到公开仓库，需按第\ref{sec:nlp-eval-contamination}小节的$n$元重叠检测删去与HumanEval、MBPP等测试集重合的代码，DeepSeek-Coder的阈值是10元组。语料中还混入自然语言，DeepSeek-Coder的配比为87\%源代码、10\%与代码相关的英文文本和3\%与代码无关的中文文本，与代码相关的文本帮助模型读懂用自然语言写的规格。

训练目标仍是下一词元预测。编辑器中的补全需要同时利用光标两侧的代码，这由第\ref{sec:nlp-llm-objective-fim}小节的中间填充（FIM）解决：把文档在字符级随机切成前缀$\bm p$、中间$\bm m$、后缀$\bm s$，用哨兵词元重排为前缀—后缀—中间（PSM）顺序；推断时以光标前、后的文件内容作为$\bm p$与$\bm s$。例如光标停在代码\ref{code:nlp-execution-sort-trace}的类型检查分支处时，模型可以读到后文的\code{sorted(set(values))}，得知输入应是可哈希的整数。用于代码模型时还要确定两点。其一是FIM比例：DeepSeek-Coder的消融中，全部样本都做FIM时填充指标最好，左到右的代码补全却最弱，最终取50\%的PSM样本。其二，FIM在打包成固定长度序列之前按文档施加，使三段来自同一个文件。

单个文件仍不足以写对代码：调用另一个模块的函数时，需要知道它的签名和返回类型。\term{仓库级上下文}（repository-level context）指同一仓库中与当前文件有依赖关系的其他文件。DeepSeek-Coder在构造预训练样本时，用正则表达式提取文件间的导入关系（Python的\code{import}、C的\code{\#include}等），按依赖拓扑排序后把同一仓库的文件拼成一个样本，每个文件前加一行路径注释；去重也改在拼接后的仓库级样本上做，以免破坏仓库结构。这样，被依赖的定义总出现在使用它的代码之前，模型在预训练中就学会了“先读依赖、再写调用”。算法\ref{alg:nlp-execution-repo-order}给出排序过程，其中按最小入度而非零入度取节点，是为了处理循环依赖。

\Needspace{12\baselineskip}
\begin{algorithm}[仓库级预训练样本的依赖排序]
  \label{alg:nlp-execution-repo-order}
  \begin{algorithmic}[1]
    \Require 一个仓库的文件集合$F$，各语言的导入语句模式
    \Ensure 若干按依赖顺序拼接的训练样本
    \State 对每个文件$b\in F$解析导入语句；若$b$引用$a$，加有向边$a\to b$，记录每个文件的入度
    \For{依赖图的每个弱连通分量$G$}
      \State $L\gets$空列表
      \While{$G$非空}
        \State 取$G$中入度最小的文件$a$（存在环时入度可能都大于零）
        \State 把$a$追加到$L$，从$G$中删去$a$并减小其后继的入度
      \EndWhile
      \State 按$L$的顺序拼接文件内容，每个文件前加路径注释，得到一个样本
    \EndFor
    \State 在拼接后的样本层面做近重复去除，再按需施加FIM与打包
  \end{algorithmic}
\end{algorithm}

推断时仓库往往远超上下文窗口，常按导入关系与第\ref{sec:nlp-execution-code-retrieval}小节的检索取出相关片段，放在当前文件之前；训练时见过“依赖在前”的格式，这种拼接才与训练分布一致。代码语言模型的收益来自数据与样本构造，代价是许可证审查、去重和去污染的治理成本。它提供候选，但单次采样的正确率仍然有限，这引出下面的筛选与修订方法。

\subsubsection{大规模采样、聚类与过滤}
\label{sec:nlp-execution-alphacode}

竞赛编程题的题面长、算法要求高，即使较强的代码模型，单次采样的正确率也很低。另一方面，pass@$k$随$k$近似按对数线性增长：多采样确实能覆盖到正确程序。问题在于实际提交次数有限，竞赛对每次错误提交都有罚时，评价协议只允许每题提交少量程序。于是任务变成：给定$n$个采样和至多$K$次提交，怎样从中挑出$K$个，使至少一个通过隐藏测试的概率尽量高。AlphaCode用记号$K$@$k$表示这一指标\cite{nlp-execution-li2022alphacode}：每题采样$k$个、只能提交其中$K$个时的解题比例，原文取$K=10$。pass@$k$相当于允许把$k$个全部提交，是$K$@$k$的上界。

AlphaCode的主要思想是以采样计算换提交正确率，分三步缩小候选集，流程见图\ref{fig:nlp-execution-alphacode}。

\begin{figure*}[htbp]
  \centering
  % 第九卷 14.2.2.2：AlphaCode 的大规模采样、示例测试过滤与按行为聚类（跨栏，约160 mm宽）。
% FigureInput：题面与测试输入；FigureModel：生成模型；FigureOutput：提交集合；中间为执行与筛选步骤。
\begin{tikzpicture}[
  x=1mm,y=1mm,
  figure picture,
  srcbox/.style={figure node,draw=FigureInput,fill=FigureInputFill,minimum height=15mm,
    minimum width=20mm,inner xsep=2pt},
  genbox/.style={figure node,draw=FigureModel,fill=FigureModelFill,minimum height=15mm,
    minimum width=29mm,inner xsep=2pt,line width=.9pt},
  procbox/.style={figure node,draw=FigureStroke,fill=white,minimum height=15mm,
    minimum width=29mm,inner xsep=2pt},
  subbox/.style={figure node,draw=FigureOutput,fill=FigureOutputFill,minimum height=15mm,
    minimum width=24mm,inner xsep=2pt}
]
  \node[srcbox] (prob) at (10,24) {题面\\示例输入输出};
  \node[genbox] (gen) at (43,24) {代码生成模型\\采样$n$个程序\\语言、标签、难度随机};
  \node[procbox] (filt) at (79,24) {示例测试过滤\\约99\%被淘汰};
  \node[procbox] (clus) at (115,24) {在新输入上执行\\按输出向量聚类};
  \node[subbox] (sub) at (150,24) {由大到小\\每类取一个\\至多$K$次提交};
  \node[srcbox,minimum width=29mm] (tin) at (115,3) {测试输入生成模型\\产生$m$个新输入};

  \draw[figure flow] (prob) -- (gen);
  \draw[figure flow] (gen) -- node[figure edge label,above=1pt] {$n$个} (filt);
  \draw[figure flow] (filt) -- node[figure edge label,above=1pt] {$n'$个} (clus);
  \draw[figure flow] (clus) -- node[figure edge label,above=1pt] {$C$类} (sub);
  \draw[figure flow] (tin) -- (clus);
  \draw[figure flow,dash pattern=on 2pt off 1.5pt] (prob.south) |- node[figure note,pos=.75,above=1pt] {以题面为条件} (tin.west);
\end{tikzpicture}

  \caption{AlphaCode的候选筛选流程。从左到右，候选数由$n$个采样经示例测试降到$n'$个，再按在新输入上的输出聚成$C$类，最后从最大的类开始每类取一个代表提交。下方的测试输入生成模型以题面为条件产生新输入，只用于区分程序行为，不需要知道正确输出。}
  \label{fig:nlp-execution-alphacode}
\end{figure*}

第一步是大规模采样，并主动制造多样性。模型是编码器—解码器结构，编码器读1536个词元的题面，解码器生成至多768个词元的程序；编码器浅、解码器深，并使用多查询注意力（所有头共享一组键值）以降低采样时的显存与带宽开销。它先在约715\,GB的GitHub代码上预训练，再在CodeContests竞赛题数据上微调。每题采样可达百万个：一半要求生成Python、一半要求生成C++，在提示中随机填入题目的算法标签和难度评分，并使用较高的采样温度。随机元数据在测试时本来就无法得知，把它们当作控制变量随机设置，增加了候选之间的差异。

第二步是用题面中的示例测试过滤。示例只有一两组输入输出，却能淘汰约99\%的采样；即便如此，许多题仍剩下成千上万个候选，约10\%的题则一个通过示例的程序都没有。第三步是按行为聚类。剩下的程序大多是写法不同、语义等价的变体，随机挑选会把有限的提交浪费在同一种行为上。若有更多测试输入，就能把行为相同的程序归为一类。设另一模型以题面为条件生成$m$个新输入$u_1,\ldots,u_m$，对每个程序$f$定义其\emph{行为签名}
\begin{equation}
 \bm b(f)=\bigl(\operatorname{run}(f,u_1),\ldots,\operatorname{run}(f,u_m)\bigr),
 \label{eq:nlp-execution-behavior-signature}
\end{equation}
其中$\operatorname{run}(f,u)$是程序在输入$u$上的输出字符串，出错或超时记为特殊标记。签名相同的程序归为一类。这里不需要知道新输入的正确输出，只比较程序之间是否一致；因此生成的输入即使不满足题目约束，只要能把不同行为分开，也有用处。提交时按类的大小从大到小，每类取一个代表；类数少于$K$时回到最大的类，跳过已提交的程序继续取。算法\ref{alg:nlp-execution-alphacode}汇总了这一选择过程。

\Needspace{12\baselineskip}
\begin{algorithm}[AlphaCode式的采样、过滤与聚类选择]
  \label{alg:nlp-execution-alphacode}
  \begin{algorithmic}[1]
    \Require 题面$x$及示例测试$E$，代码模型，测试输入生成模型，采样数$n$，提交上限$K$
    \Ensure 至多$K$个待提交程序
    \State 采样$n$个程序：语言、标签与难度在提示中随机设置
    \State $S\gets$通过$E$中全部示例的程序
    \State 由测试输入生成模型以$x$为条件生成输入$u_1,\ldots,u_m$
    \State 对每个$f\in S$执行得到签名$\bm b(f)$（式\eqref{eq:nlp-execution-behavior-signature}），按签名相等划分为类，类按大小降序排列
    \State $R\gets$空列表
    \While{$|R|<K$且$S$中仍有未提交的程序}
      \For{按大小降序的每一类}
        \If{该类中有未提交的程序且$|R|<K$}
          \State 取一个未提交的程序加入$R$
        \EndIf
      \EndFor
    \EndWhile
    \State \Return $R$
  \end{algorithmic}
\end{algorithm}

为什么优先取大类？原文的解释是，错误程序可以错得各不相同，正确程序在任何输入上都给出相同输出，因而容易汇聚成大类。一个教学示例可以说明这一点及其边界。题目要求输出整数列表去重后的升序结果，示例为\code{[3,1,3]}$\to$\code{[1,3]}。设100个采样中29个通过示例，生成的新输入为\code{[]}与\code{[2,-1,2,5]}。其中18个正确，7个在空输入上出错，4个错误地按绝对值排序。按绝对值排序在这两个新输入上恰好也输出\code{[]}和\code{[-1,2,5]}，于是与正确程序同签名，类大小为22、7。取$K=2$时，第一次提交从22个程序中取一个，正确的概率为$18/22$；第二次取空输入出错的那一类，不会再补救。若新输入中有\code{[-3,1]}，按绝对值排序会输出\code{[1,-3]}，两类就被分开。这说明聚类的效果取决于生成输入区分行为的能力。

与其他方法相比，按签名聚类后取大类，相当于在程序行为上做多数投票：第\ref{sec:nlp-execution-self-consistency}小节的自一致性对最终答案字符串投票，这里对整组输出投票，所比较的对象从一个答案变成了一个函数在多个输入上的取值。它的优点是不需要训练额外的评分器，执行结果也比模型打分更客观；代价是采样与执行的计算量巨大，且依赖示例测试和生成输入的质量。原文的对照显示，过滤加聚类使解题率随采样数持续增长，不过滤时几乎不随采样数增长，但与“允许提交全部候选”的上界之间仍有很大差距。

\subsubsection{语法约束解码：SynCode}
\label{sec:nlp-execution-syncode}

即使规格清楚，自由生成仍可能漏括号、缩进错误或在JSON字段之间遗漏逗号，语法错误的程序连测试都无法运行。第\ref{sec:nlp-inference-grammar}小节已经说明文法约束解码的一般机制：在每一步算出可以接续当前前缀的合法词元集合，把其余词元的对数分数设为$-\infty$后再选择。SynCode是面向编程语言文法的一种实现\scite{nlp-execution-ugare2024syncode}，这里只讲它怎样处理模型词元与文法终结符不对齐的问题。

编程语言的文法分两层。词法层用正则表达式定义\term{终结符}（terminal），例如标识符、整数、字符串、关键字和各种括号，每个终结符对应一个确定有限自动机（DFA）；语法层用上下文无关文法规定终结符怎样组合。模型的词元却是按频率切出的子词，与终结符没有对应关系：一个词元可能只是某个标识符的一半，也可能跨越两个甚至三个终结符。例如已生成的JSON前缀为
\begin{quote}
\code{\{"date":"2025-04-}
\end{quote}
一个词元可能解码为\code{18","}，同时补完日期字符串、产生逗号并开启下一个字段名的字符串。只知道“下一个终结符是字符串”不足以判断这个词元是否合法。

SynCode把判断拆成语法层的在线计算和词法层的离线查表。在线部分先对当前前缀做词法分析，把它分成已经确定的终结符序列和\emph{剩余后缀}$r$：若前缀以一个尚未闭合的字符串结束，$r$就是这段未完成的字符串\code{"2025-04-}；若前缀以一个完整的词法单元结束，$r$取这个单元，因为它可能被继续延长而改变类型，例如\code{retu}之后可能变成关键字\code{return}。确定的终结符序列交给增量的LR(1)解析器，解析器状态直接给出当前位置和下一位置可接受的终结符。由此构造\term{接受序列}（accept sequence）集合$\mathcal A$：每个接受序列$\Lambda=(\tau_1,\ldots,\tau_\alpha)$是一串终结符，把它接在已确定部分之后仍是合法前缀，其中$\tau_1$是$r$所属或将要成为的终结符。实现中只用长度为1和2的接受序列。

离线部分是\term{DFA掩码表}（DFA mask store）$\mathcal M_{\alpha}$。对终结符DFA中的每个状态$q$和后续终结符序列$\Lambda'$（长度为$\alpha$），表项$\mathcal M_{\alpha}(q,\Lambda')$是一个长为词表大小$|V|$的布尔向量。词元$t$的位置为真，当且仅当从状态$q$读入$t$的字符串时满足以下之一：读完后仍停在当前终结符DFA的活状态（可以继续扩成合法单元）；或者在中途某处到达当前终结符的接受状态，剩余字符能被$\Lambda'$中的下一个终结符继续匹配，依此递归；或者$\Lambda'$已用完，剩余字符不再检查。最后一种情形保证了不会误删合法词元，也是不完备的来源。生成第$k+1$个词元时，掩码为
\begin{equation}
 \bm m=\bigvee_{\Lambda\in\mathcal A,\ q_r\ \text{为活状态}}\mathcal M_{|\Lambda|-1}\bigl(q_r,\Lambda_{2:}\bigr),
 \qquad q_r=\delta^*_{\tau_1}\bigl(q_0^{\tau_1},r\bigr),
 \label{eq:nlp-execution-syncode-mask}
\end{equation}
其中$\delta^*_{\tau_1}$是终结符$\tau_1$的DFA在读入字符串后的转移，$q_0^{\tau_1}$是它的初始状态，$\Lambda_{2:}$是去掉第一个终结符后的序列，$\vee$为按位或。式\eqref{eq:nlp-execution-syncode-mask}的含义是：先让剩余后缀在所属终结符的DFA中走到状态$q_r$，再按每条可能的终结符接续查表，把所有可行接续允许的词元合并。在线代价只是解析器查询和若干次布尔向量合并，词表上的逐词元字符串匹配都在离线完成。算法\ref{alg:nlp-execution-syncode}给出一步解码的过程。

\Needspace{12\baselineskip}
\begin{algorithm}[SynCode的一步约束解码]
  \label{alg:nlp-execution-syncode}
  \begin{algorithmic}[1]
    \Require 当前前缀$C_k$，LR(1)解析表，各终结符的DFA，离线构造的$\mathcal M_0,\mathcal M_1$，模型对数分数$\bm z\in\mathbb R^{|V|}$
    \Ensure 下一词元$t$
    \State 对$C_k$做词法分析，得到已确定的终结符序列与剩余后缀$r$
    \State 从缓存的解析器状态出发，增量解析新确定的终结符，得到当前与下一位置可接受的终结符
    \State 构造长度为1或2的接受序列集合$\mathcal A$
    \State $\bm m\gets\bm 0$
    \For{$\Lambda\in\mathcal A$}
      \State $q_r\gets\delta^*_{\tau_1}(q_0^{\tau_1},r)$
      \If{$q_r$是活状态}
        \State $\bm m\gets\bm m\vee\mathcal M_{|\Lambda|-1}(q_r,\Lambda_{2:})$
      \EndIf
    \EndFor
    \State 对$m_v=0$的词元令$z_v\gets-\infty$，按第\ref{chap:nlp-model-inference}章的选择规则得到$t$
    \State 缓存解析器状态，令$C_{k+1}\gets C_kt$
  \end{algorithmic}
\end{algorithm}

回到JSON例子。剩余后缀\code{"2025-04-}使字符串DFA停在“字符串内部”状态；若接受序列为（字符串，逗号），查表时\code{18","}先在\code{18"}处闭合字符串，再由逗号匹配\code{,}，余下的\code{"}因序列已用完不再检查，于是被放行；\code{18"]}则在闭合字符串后遇到\code{]}，与逗号或右花括号都不匹配而被屏蔽。设$A(C_k)$为真正能扩成合法程序的下一词元集合，$M(C_k)$为掩码放行的集合。SynCode证明了在接受序列取到全部长度为$d$的序列时$A(C_k)\subseteq M(C_k)$，即不误删合法词元（原文称可靠性）；当$d$大于所有词元的字符长度时反向包含也成立（完备性）。实现只用长度1和2，所以是可靠而不完备的：偶尔放行的非法词元会在下一步的解析中暴露。

与第\ref{sec:nlp-inference-grammar}小节XGrammar那类直接在字符级下推自动机上判断词元的做法相比，SynCode把上下文无关部分交给LR解析器、把词法部分预编译为DFA状态上的掩码，适合文法已写成LR(1)形式的编程语言。离线构造代价与DFA状态数、词表大小和终结符数成正比，原文对每一组文法与模型可在数分钟内完成。它的局限也来自分层：文法须是无歧义的LR(1)或LALR(1)文法，Python缩进这类不能由上下文无关文法表达的部分，要靠额外的缩进跟踪机制屏蔽违规词元；更根本的是，语法正确不判断语义，\code{"2025-02-30"}是合法字符串而不是有效日期，函数能被解析也可能在测试中失败。

\subsubsection{执行反馈修订：Self-Debugging}
\label{sec:nlp-execution-self-debugging}

候选能被解析以后，测试才能揭示行为错误。AlphaCode丢弃失败候选并继续采样，但失败候选往往只差一处比较或一个边界条件，失败信息本身指出了该改哪里。Self-Debugging把执行结果写回提示，让同一个模型看到自己的代码和反馈后修订\scite{nlp-analysis-chen2024selfdebug}。它用少样本提示实现，示例演示“代码—反馈—解释—修订”的过程，模型参数不变；原文把其中的代码解释类比为\term{橡皮鸭调试}（rubber duck debugging），即程序员向他人（哪怕是一只橡皮鸭）逐行讲解代码时，常会自己发现错误。

原文比较了四种反馈格式，信息量依次增加：简单反馈只说“上面的代码不正确，请修正”；单元测试反馈给出失败用例的调用、预期输出与实际输出或异常；代码解释反馈在此基础上要求模型先逐行解释自己的代码，再判断与任务描述哪里不符；执行轨迹反馈则要求模型逐行描述代码在失败输入上的执行过程。后两种格式中的解释和轨迹描述都由模型生成，只有测试结果来自真实执行。在没有单元测试的任务（如原文中的文本到SQL）上，修订的依据完全来自模型自己的解释。

以本节的截止判断为例，设候选把“4月16日18:00前提交”实现为\code{<=}。代码\ref{code:nlp-execution-feedback}是一轮单元测试加解释的反馈格式。

\begin{codeblock}[label={code:nlp-execution-feedback}]{一轮Self-Debugging的提示内容（教学示例）}
[任务] 校外来宾须在截止时刻之前提交；恰好在截止时刻提交应拒绝。
[代码]
def eligible(kind, submitted_at, deadline):
    ...
    return to_time(submitted_at) <= to_time(deadline)
[执行反馈]  # 由测试运行器生成
test_exact_deadline 失败:
  eligible("external", "2025-04-16T18:00:00+08:00",
           "2025-04-16T18:00:00+08:00")
  预期 False, 实际 True
[解释]  # 由模型生成
最后一行在两个时刻相等时返回 True, 而任务要求“之前”, 即严格小于。
[修订后的代码]
\end{codeblock}

模型据此把比较符改为\code{<}，再运行同一失败用例和17:59允许、次日拒绝等回归用例。反馈应包含输入、预期、实际返回和测试来源；“边界不对”过于含糊，“请改成小于”则已直接给出补丁，二者都不适合用来比较自主修订能力。若时间解析器此前已把时刻截成日期，只改比较符会把17:59也拒绝，下一轮反馈会要求继续追踪调用链（第\ref{sec:nlp-execution-repair}小节）。算法\ref{alg:nlp-execution-self-debugging}给出轮次结构。

\Needspace{12\baselineskip}
\begin{algorithm}[Self-Debugging的多轮修订]
  \label{alg:nlp-execution-self-debugging}
  \begin{algorithmic}[1]
    \Require 规格$x$，开发测试$T$（可为空），带调试示例的少样本提示，反馈格式，轮数上限$R$
    \Ensure 候选程序及其状态
    \State $y^{(0)}\gets$模型以$x$为条件生成的程序
    \For{$i=0,1,\ldots,R-1$}
      \State 在$T$上运行$y^{(i)}$，记录每个用例的输入、预期、实际输出或异常
      \If{$T$非空且全部通过}
        \State \Return $y^{(i)}$，状态为通过开发测试
      \EndIf
      \State 按反馈格式组装$f^{(i)}$：失败用例；可选地请模型逐行解释$y^{(i)}$或其在失败输入上的执行过程
      \If{$T$为空且模型的解释判定$y^{(i)}$正确}
        \State \Return $y^{(i)}$，状态为模型自评正确
      \EndIf
      \State $y^{(i+1)}\gets$模型以$(x,y^{(i)},f^{(i)})$为条件生成的修订
      \If{$y^{(i+1)}=y^{(i)}$}
        \State \textbf{break}\Comment{修订不再改变代码}
      \EndIf
    \EndFor
    \State \Return 最后一个候选，状态为未解决
  \end{algorithmic}
\end{algorithm}

Self-Debugging与AlphaCode代表了利用执行信息的两种方式：后者并行采样、只用测试做筛选，前者串行修订、把测试结果当作改写的条件；二者可以组合，例如先采样若干候选，再对失败者分别修订。比较两者时应固定总的模型调用、生成词元和测试执行预算，否则多轮的收益会被误归因于反馈机制。Self-Debugging对失败候选的信息利用更充分；局限是修订受限于模型理解反馈的能力，没有测试时模型自评“现在正确”不增加独立证据。

\subsection{代码转换}
\label{sec:nlp-execution-transformation}

转换任务的输入多了一份原程序，目标是把它迁到另一种语言或接口，同时保持约定的行为。难点首先在于“保持行为”不能从文本外观判断。对$-3$和2，C++与Java的整数除法向零截断得到$-1$，Python的\code{-3 // 2}向下取整得到$-2$，用\code{/}则得到浮点数$-1.5$；固定宽度无符号32位加法把$2^{32}-1+1$回绕为0，Python整数则得到$2^{32}$，需要按目标契约决定是否显式取模。异常类型、时区库和可变对象的别名也可能改变行为。所以转换后要在同一组输入上比较返回、异常和副作用，结论仅限于所测输入范围。

方法上有两条路线。基于语法树的规则转换精确、可审查，但规则须人工编写；学习式翻译能处理任意构造，但平行代码（同一功能在两种语言下的实现）稀少。回译式无监督翻译则只用两种语言各自的单语代码训练。

\subsubsection{AST规则转换}
\label{sec:nlp-execution-ast-rewrite}

接口迁移和版本升级往往只需改动少数有固定模式的构造，例如把记录键\code{published\_at}改为新接口的\code{publishedAt}，同时保持按时刻排序的行为。规则转换先限定原代码与可改的接口，再在AST上按模式匹配节点并重写：这里只匹配目标函数中以该字符串常量为下标的访问节点，替换常量后重新生成源码；全局文本替换则可能误改文档示例或另一个业务对象中的同名字段。Python可以用\code{ast.NodeTransformer}按条件替换下标节点，再用\code{ast.unparse}生成代码。

转换后要检查语法是否正确、目标键是否替换完整、允许范围外的节点是否保持原样。规则转换的知识完全来自人工编写的匹配与重写条件，行为可预测、可逐条审查；代价是每类迁移都要写规则，规则覆盖之外的构造（例如通过变量间接传入的键名）需人工处理。跨语言的整体迁移涉及大量语义差异，手写规则的转译器（transcompiler）成本高，产生的代码也常不符合目标语言习惯，这引出学习式翻译。

\subsubsection{回译式无监督代码翻译}
\label{sec:nlp-execution-transcoder}

学习式翻译若用平行数据监督训练，需要大量“同一函数的C++版与Python版”，这类数据很少。TransCoder把第\ref{sec:nlp-generation-unsupervised-mt}小节的无监督机器翻译搬到代码上，在C++、Java和Python之间翻译函数，只用GitHub上的单语代码训练\scite{nlp-execution-lachaux2020transcoder}。无监督翻译的三个要素在代码上各有对应的实现。

模型是一个在三种语言间共享的编码器—解码器Transformer，词表是在三种语言代码上联合学习的子词词表，每种语言先用各自的词法分析器切分，使字符串和注释的处理符合该语言规则。跨语言初始化采用跨语言掩码语言模型预训练：把三种语言的源代码混在一起，随机遮蔽词元并预测，用得到的编码器初始化翻译模型的编码器与解码器。代码之所以能在没有平行数据的情况下对齐，是因为三种语言共享大量锚点：\code{for}、\code{while}、\code{if}等关键字，数字、运算符，以及\code{min}、\code{max}这类同名库函数，它们在不同语言中出现的上下文相近，把周围的词元拉到同一表示空间。

第二个要素是去噪自编码，对应式\eqref{eq:nlp-generation-unsup-dae}。加噪函数随机遮蔽、删除词元并在小窗口内打乱顺序，模型须恢复原函数；解码器开头的语言标记指定输出语言。单纯复制无法完成这个任务，解码器因此学会生成合乎该语言语法的代码，对输入噪声也更稳健。第三个要素是回译，对应式\eqref{eq:nlp-generation-unsup-bt}：用当前模型把Python函数译成C++，再训练模型从这份可能不完美的C++恢复原Python函数，反之亦然。随着训练进行，正向翻译变好，产生的伪平行对质量也提高，二者相互促进。

评价上，TransCoder的作者认为BLEU或参考匹配不足以衡量代码翻译：变量改名或等价写法会被判错，一个比较符的错误却只扣很少的分数。他们从GeeksforGeeks网站收集了852个有三种语言实现的函数，配上单元测试，以\term{计算准确率}（computational accuracy）评价，即翻译结果在相同输入上与参考实现输出一致的比例。束搜索给出多个译文时，至少一个通过测试即记为正确，这与pass@$k$的思路相同。

回译式翻译与规则转换互补：它不需要写规则，能处理库调用的映射和惯用写法；但翻译只受语言模型与回译的统计约束，不保证行为一致，本小节开头的整数除法就是典型错误。学习式翻译的结果因此须经单元测试核验，负数、溢出和异常等边界输入应写进测试。

\subsection{仓库级程序修复}
\label{sec:nlp-execution-repair}

仓库修复以既有仓库和问题描述为输入，交付一个补丁，与从零生成的区别在于大部分行为必须保持不变。SWE-bench从真实的GitHub问题及其修复提交中构造任务：模型看到问题描述和修复前的仓库版本，生成的补丁要使修复提交附带的失败测试转为通过，同时使原本通过的测试保持通过\scite{nlp-intro-jimenez2024}。挑战有三个：在成百上千个文件中定位应修改的位置；正确编辑跨文件的代码，不引入语法或缩进错误；在没有标准答案的情况下验证补丁。下面先介绍让模型自主完成这三件事的SWE-agent，再介绍改用固定流程的Agentless，最后说明多候选补丁怎样借测试选择。

本小节以一个两文件教学仓库作为共同基线。按其业务接口，校内人员直接返回允许，函数\code{eligible}只判断校外来宾提交信息是否早于截止时刻；时间字符串必须带时区，比较前统一到北京时间。共同通知发布于2025年4月15日，活动\code{ML-0418}改至4月18日周五14:00，校外提交截止为4月16日18:00。问题报告为“截止当刻仍被允许，同一天的通知排序也不正确”。相关代码如下。

\begin{codeblock}[language=Python,label={code:nlp-execution-repair-base}]{两个文件中相互关联的日期错误}
# time_utils.py
from datetime import datetime, timedelta, timezone
LOCAL = timezone(timedelta(hours=8))

def parse_local(value):
    timestamp = datetime.fromisoformat(value)
    if timestamp.tzinfo is None:
        raise ValueError("timezone required")
    return timestamp.astimezone(LOCAL).date()

# notices.py
from time_utils import parse_local

def eligible(kind, submitted_at, deadline):
    if kind == "internal":
        return True
    if kind != "external":
        raise ValueError("unknown participant kind")
    return parse_local(submitted_at) <= parse_local(deadline)

def sort_notices(rows):
    def key(row):
        return parse_local(row["published_at"])
    return sorted(rows, key=key)
\end{codeblock}

这里有两处独立错误：\code{.date()}使09:00与14:00都变成4月15日，排序只能保留输入次序；\code{<=}即使收到完整时刻，也会接纳恰好18:00的提交。正确修订需要去掉日期截断，并把报名比较改成\code{<}，排序函数本身无需修改。

\subsubsection{智能体–计算机接口：SWE-agent}
\label{sec:nlp-execution-swe-agent}

让语言模型修仓库，最直接的办法是给它一个命令行，在第\ref{sec:nlp-execution-react}小节的ReAct式循环中交替写出思考与命令，读取命令输出后继续。问题在于命令行是为人设计的：\code{cat}一个两千行的文件会把上下文塞满；用\code{sed}编辑时行号或缩进稍有差错就悄悄写坏文件；命令成功时常常什么也不输出，模型无法确认操作是否生效。人类程序员靠编辑器和集成开发环境克服这些不便，SWE-agent的观点是，语言模型也是一类新的用户，同样值得为它专门设计界面\scite{nlp-execution-yang2024sweagent}。

\term{智能体–计算机接口}（agent-computer interface, ACI）指智能体与计算机之间的一组命令及其反馈格式。SWE-agent的接口遵循几条设计原则：命令简单、少而易懂，带简洁的文档和示例；常用操作（如查看、搜索、编辑）合并为一条命令完成，减少多步操作；反馈信息简洁而完整，让模型知道环境当前状态和刚才操作的效果；内置防护，在错误扩散之前拦截。表\ref{tab:nlp-execution-aci}列出主要命令。

\begin{table}[htbp]
  \centering
  \small
  \begin{tabularx}{\linewidth}{@{}>{\raggedright\arraybackslash}p{0.3\linewidth}>{\raggedright\arraybackslash}X@{}}
    \toprule
    命令 & 作用与反馈 \\
    \midrule
    \code{find\_file}、\code{search\_dir}、\code{search\_file} & 按文件名、目录内或文件内搜索字符串；返回匹配的文件与行号摘要，超过50条时不返回结果而提示缩小查询 \\
    \code{open}、\code{goto}、\code{scroll\_up}、\code{scroll\_down} & 文件查看器，每次显示100行并带行号，注明文件路径、总行数及窗口上下省略的行数 \\
    \code{edit}、\code{create} & 替换当前文件中指定行区间的内容或新建文件；编辑后自动运行代码检查器，若引入语法、缩进或未定义名称错误则撤销编辑，返回错误类型及编辑前后的片段 \\
    \code{submit} & 以当前工作区相对基线版本的差分作为补丁提交 \\
    \bottomrule
  \end{tabularx}
  \caption{SWE-agent接口的主要命令。普通命令行命令（如运行测试脚本）仍可使用；命令成功而无输出时，接口明确告知这一点。}
  \label{tab:nlp-execution-aci}
\end{table}
\FloatBarrier

三类命令分别针对定位、查看和编辑中的典型失误。搜索命令只返回摘要并限制条数；查看器的窗口大小经过消融选择，原文的对照中100行窗口的解题率高于30行窗口和整文件显示：30行窗口要反复滚动才能看全一个函数，整文件显示则带入大量无关内容。编辑命令与\term{代码检查器}（linter）绑定，在写坏文件之前拦截错误，模型据此重试，不必等到测试失败后才回头查找原因。上下文管理也是接口的一部分：只保留最近5次观察的完整内容，更早的观察折叠为一行占位。代码\ref{code:nlp-execution-aci-trace}是在教学仓库上的一段示意轨迹。

\begin{codeblock}[label={code:nlp-execution-aci-trace}]{SWE-agent式轨迹片段（教学示意）}
思考: 问题涉及截止时刻, 先找 eligible 的定义。
动作: search_dir "def eligible"
观察: Found 1 matches for "def eligible" in /repo:
      /repo/notices.py (1 matches)
思考: 打开该文件查看比较逻辑。
动作: open notices.py
观察: [File: /repo/notices.py (13 lines total)]
      1:from time_utils import parse_local
      ...
      8:    return parse_local(submitted_at) <= parse_local(deadline)
思考: parse_local 的返回类型决定比较粒度, 先看它的定义。
动作: open time_utils.py
观察: ...
      8:    return timestamp.astimezone(LOCAL).date()
思考: .date() 丢掉了时刻; 去掉它, 再把比较改为严格小于。
动作: edit 8:8
          return timestamp.astimezone(LOCAL)
      end_of_edit
观察: File updated. ...
\end{codeblock}

算法\ref{alg:nlp-execution-swe-agent}概括了整个循环。模型每一步输出一段思考和一条命令；格式不合法时，接口返回错误信息让模型重写。循环在模型调用\code{submit}或耗尽预算时结束。

\Needspace{12\baselineskip}
\begin{algorithm}[SWE-agent的交互循环]
  \label{alg:nlp-execution-swe-agent}
  \begin{algorithmic}[1]
    \Require 基线版本的仓库与执行环境，问题描述，接口命令及其文档，完整保留的观察数$h=5$，调用预算
    \Ensure 补丁或未完成状态
    \State 历史$H\gets$（系统提示：命令文档与一条示范轨迹；问题描述）
    \While{预算未耗尽}
      \State 模型以$H$为条件输出（思考，命令）
      \If{输出格式不合法}
        \State 把格式错误信息作为观察，要求重写；\textbf{continue}
      \EndIf
      \If{命令为\code{submit}}
        \State \Return 工作区相对基线的差分
      \EndIf
      \State 在环境中执行命令，按接口格式化观察（查看窗口、搜索摘要、检查器结果，或“成功且无输出”）
      \State 把（思考，命令，观察）追加到$H$；把最近$h$次之前的观察折叠为一行
    \EndWhile
    \State \Return 未完成，附当前差分与轨迹
  \end{algorithmic}
\end{algorithm}

SWE-agent说明接口设计本身就是提升智能体能力的一个维度：同一模型、同一循环，只改变命令与反馈格式，解题率就有明显差别。它的优点是通用，模型可以按需复现问题、运行测试、多次尝试；代价是轨迹长、调用次数多，早期的错误判断（如打开了错误文件）会沿轨迹累积，整个流程依赖模型自己决定下一步做什么。

\subsubsection{分层定位与局部修复：Agentless}
\label{sec:nlp-execution-agentless}

Agentless针对的正是自主智能体的这些问题\scite{nlp-execution-xia2024agentless}。它的作者指出，当前模型的规划与自我纠错能力有限：一旦走错方向，长轨迹中的后续动作难以回到正轨；每一步都要读写上下文，成本高；轨迹又长又多样，出错后难以分析原因。Agentless不让模型做流程决策，把修复拆成定位、修复、补丁验证三个固定阶段，模型只在每个阶段内部完成定义明确的子任务。

定位分三层逐步缩小。第一层，把仓库的目录结构交给模型，请它列出最可疑的文件，并用嵌入检索补充文件候选；第二层，只给出这些文件的骨架，即类、函数和全局变量的声明，请模型选出相关元素；第三层，读取这些元素的完整代码，确定具体的编辑位置。分层的好处是每一层的输入都较短，又不会过早丢弃跨文件的线索。在教学仓库上，目录与问题描述指向\code{notices.py}，骨架中的\code{eligible}和\code{sort\_notices}都调用\code{parse\_local}，于是第三层定位到\code{time\_utils.py}的返回语句和\code{notices.py}的比较。若定位器只提供\code{eligible}，修复器就看不到辅助函数的返回类型，容易写出下一小节中的回归补丁。

修复阶段给出编辑位置附近的上下文，请模型以“搜索—替换”的差分格式写出局部补丁，即指明要被替换的原代码片段和替换后的代码，并多次采样得到多个候选。差分格式比重写整个文件短，也比按行号编辑更不容易错位。补丁验证阶段用测试筛选候选：模型先生成一个复现问题的测试，并在原仓库上运行确认它能复现失败；随后对每个候选运行复现测试和仓库原有的回归测试，再对规范化后的补丁做多数投票。

与SWE-agent相比，Agentless以灵活性换取可控性：流程固定，每一步的输入输出都可检查，调用次数和成本可以预先估计；代价是流程之外的情形无法自行调整，例如问题需要先安装依赖、运行脚本才能定位时，固定流程就没有对应步骤。两者都离不开多候选补丁与测试选择，这是下一小节的内容。

\subsubsection{多候选补丁与测试选择}
\label{sec:nlp-execution-patch-selection}

在相同基线上构造三个教学候选：$P_1$只保留完整时刻，$P_2$只改严格比较，$P_3$同时修正两处。目标测试有两项：截止当刻应拒绝，同日通知应按时刻排序。七项回归分别检查截止前允许、次日拒绝、校内直接允许、空列表、重复项与稳定次序、缺时区报错以及未知人员类型报错。每个候选在独立子进程中导入自己的两文件副本（Python 3.9.6，5秒超时），避免模块缓存把上一候选带入下一候选。表\ref{tab:nlp-execution-patches}记录实际运行结果。

\begin{table}[htbp]
  \centering
  \begin{tabular}{lcc}
    \toprule
    候选 & 目标测试通过 & 回归通过 \\
    \midrule
    原程序 & 0/2 & 7/7 \\
    $P_1$：仅保留时刻 & 1/2 & 7/7 \\
    $P_2$：仅严格比较 & 1/2 & 6/7 \\
    $P_3$：两处共同修正 & 2/2 & 7/7 \\
    \bottomrule
  \end{tabular}
  \caption{跨文件日期修复的教学运行结果。分母固定，失败均为行为断言不符，不是导入或依赖故障。}
  \label{tab:nlp-execution-patches}
\end{table}
\FloatBarrier

$P_2$通过了截止当刻拒绝，却因日期截断把17:59也变成与截止同一天的日期，严格小于为假，回归测试“截止前允许”失败。可见“原失败测试转为通过”不足以选择补丁。教学选择器要求目标测试全部通过、约定回归全部保持，并检查差分只修改允许的位置，因而选出$P_3$；多个候选都通过时，再按预定的修改范围与平局规则选择。把$P_1$收到的边界失败反馈接入第\ref{sec:nlp-execution-self-debugging}小节的修订也会得到$P_3$，但顺序修订与独立多候选采样应作为不同协议记录。

冻结$P_3$后另运行四个预先留出的验收用例：等价UTC偏移的截止时刻、截止前一微秒、无效日历日期和无效记录类型，四项全部通过。Agentless原文的选择规则与教学选择器不同：先保留回归失败最少的候选，再用复现测试过滤，最后做多数投票，复现测试不可靠时有回退策略。无论采用哪种规则，模型生成的测试的预期都应对照问题独立审查；还应限制补丁可修改的文件范围、拒绝删改测试，以免候选通过改变验收条件来消除失败。

\section{工具与环境任务}
\label{sec:nlp-execution-environment}

程序任务的产物可以在固定副本上反复运行测试，环境中的一次修改却会留下来，影响后续所有操作。用户说“把我的预约改到通知中的时间和地点，不要发送通知，保留备注”，系统需要找到正确的记录，读取当前状态，执行允许的修改，再核验结果。工具返回、用户补充以及其他入口的修改都可能改变下一步该做什么，因此这类任务的基本单位不再是一次完整的文本生成，而是带观察的行动。本节把这样的系统称为\term{智能体}（agent）：语言模型在循环中依据目标和已有观察选择下一步，或调用工具，或回复用户；外部执行器负责真正执行，并把结果作为观察返回。

这一设定带来一串相互衔接的问题：模型先要获得调用工具的能力；工具很多时，要选出合适的工具并把参数绑定到可靠来源；多步任务要安排行动顺序，在观察出乎意料时调整，并让独立的子任务并行；轨迹变长后，要保存仍然有效的状态、取回相关历史，并把失败经验带到下一次尝试；提示驱动的能力不足时，要用交互轨迹训练模型；单个上下文难以兼顾执行与检查时，可以让多个模型实例分工；最后，执行失败要能恢复，完成与否要由环境证据判断。表\ref{tab:nlp-execution-agent-map}按这些挑战列出代表方法：前一行方法留下的缺口，往往就是后一行方法的出发点。

\begin{table*}[htbp]
\centering
\small
\begin{tabular}{p{18mm}p{30mm}p{48mm}p{44mm}}
\toprule
挑战 & 代表方法 & 核心机制 & 新增代价或局限 \\
\midrule
获得调用能力 & 函数调用格式、Toolformer、ToolLLM、Gorilla & 工具描述进入上下文，调用写成带特殊标记的结构化片段；自监督插入并按损失下降过滤；强模型合成指令与调用轨迹 & 合成数据有覆盖与正确性偏差；自监督只适合无副作用、可批量执行的工具 \\
选工具与填参数 & 描述嵌入检索、模式约束解码 & 按缺失信息检索候选工具，再查权限与副作用；参数逐项绑定来源，模式保证类型合法 & 合法的值仍可能指向错误对象 \\
规划 & Plan-and-Solve、ReWOO、ReAct、LLMCompiler & 先写计划再执行；每步依据最新观察决定；按依赖图并行派发 & 先规划难以应对意外；逐步交替调用多、上下文增长；并行依赖正确声明依赖 \\
记忆 & 结构化状态、Generative Agents、MemGPT & 按事件更新带来源与版本的状态；按新近度、重要性、相关性打分取回；分层上下文由模型自行换入换出 & 摘要可能丢失否定与未知项；检索命中相似但不适用的记录 \\
反思 & Reflexion & 把失败轨迹与评价写成文字教训，作为下一回合的上下文 & 不改参数；反思可能错误归因 \\
训练 & FireAct、AgentTuning；多轮强化学习（Search-R1、ArCHer） & 成功轨迹微调并混合通用数据；整条轨迹以任务结果为奖励，观察词元不计梯度 & 只能模仿示范；奖励稀疏、信用分配难、环境重置昂贵 \\
协作 & 角色分工（CAMEL、MetaGPT）、多智能体辩论 & 不同角色处理规划、执行与审查；多个实例互读答案后修订并聚合 & 通信与调用成本成倍增加；错误可在实例间传播 \\
恢复与完成 & 幂等键、独立回读、终态判定、pass\textasciicircum{}$k$ & 按失败类型分别恢复；完成以环境读取为准；重复运行衡量稳定性 & 依赖服务端契约；环境不能无损重置时重复试验不独立 \\
\bottomrule
\end{tabular}
\caption{工具与环境任务的挑战与方法谱系。每行列出一类挑战的代表方法、它们改变的机制以及随之带来的新代价；后续各小节按表中顺序展开。}
\label{tab:nlp-execution-agent-map}
\end{table*}

为了让各方法在同一任务上比较，本节使用一组教学业务状态。用户\code{U7}有个人预约\code{R17}，关联活动\code{ML-0418}，原时间为2025年4月17日14:00，地点字段为海棠楼101，备注为“带笔记”，版本为1。同一用户另有同名预约\code{R18}，关联其他活动，必须保持不变。共同通知给出的目标为4月18日周五14:00、海棠楼201；修改个人预约不修改活动通知，也不办理报名。模拟环境以\code{Haitang-201}表示房间，所有时间带北京时间偏移。配套程序是运行在本地内存中的确定性模拟，没有调用真实预约系统；以下轨迹、调用数与数值均为教学构造，不代表任何方法的实测效果。

\subsection{函数调用能力的获得}
\label{sec:nlp-execution-function-calling}

模型要使用工具，须先知道有哪些工具、调用应写成什么格式，以及在什么时候调用。最直接的做法是把工具描述写进提示，再用含调用的示范做有监督微调，但示范要人工编写，覆盖的工具有限。Toolformer让模型在普通文本上自己尝试插入调用，只保留确实有助于预测后文的那些；ToolLLM与Gorilla用强模型大规模合成指令与调用轨迹，把覆盖范围扩大到成千上万个真实接口。三条路线都要先约定调用在词元序列中的表示。

\subsubsection{工具描述与调用格式}
\label{sec:nlp-execution-call-format}

函数模式的写法（函数名、用途、参数的JSON模式）以及执行器校验参数、检查权限、注入身份的“调用—结果—回复”循环，已在第\ref{sec:nlp-dialogue-llm-function}小节给出。这里只补模型一侧：一次调用在词元序列中长什么样，模型怎样学会写出它。代码\ref{code:nlp-execution-update-schema}是教学目录中更新工具的描述，另附返回值说明；用途说明中的否定条件“不发送任何通知”决定了它能否用于本例。

\begin{codeblock}[label={code:nlp-execution-update-schema}]{更新工具的描述（教学示例）}
{
  "name": "update",
  "description": "修改单条个人预约的开始时间和房间；不发送任何通知，不改备注。",
  "parameters": {
    "type": "object",
    "properties": {
      "id":               {"type": "string"},
      "expected_version": {"type": "integer", "minimum": 1},
      "start":            {"type": "string", "format": "date-time"},
      "room":             {"type": "string"},
      "request_id":       {"type": "string"}
    },
    "required": ["id", "expected_version", "start", "room", "request_id"],
    "additionalProperties": false
  },
  "returns": "回执：status、new_version；版本不符时返回 CONFLICT"
}
\end{codeblock}

模型看到和写出的都是词元序列，因此工具描述、调用与结果必须按固定的\term{对话模板}（chat template）序列化。常见做法是：系统消息中列出全部工具描述；模型要调用工具时，在助手消息中输出一段由特殊标记包裹的结构化调用；执行器把结果写成一条工具消息追加到序列末尾，模型再继续生成。代码\ref{code:nlp-execution-call-template}给出一种教学模板，其中\code{<|assistant|>}、\code{<tool\_call>}等是加入词表的特殊词元。不同模型的模板写法不同，有的用标签包裹JSON，有的直接生成类似Python的函数调用；训练与推断必须使用同一套模板。

\begin{codeblock}[label={code:nlp-execution-call-template}]{含并行调用的对话序列化（教学模板）}
<|system|>可用工具：[{"name":"search",...},{"name":"read",...},
  {"name":"availability",...},{"name":"update",...}]<|end|>
<|user|>把我的读书会预约改到通知中的时间和地点，不要通知，保留备注。<|end|>
<|assistant|><tool_call>{"id":"c1","name":"search",
  "arguments":{"user":"U7","event":"ML-0418"}}</tool_call><|end|>
<|tool|>{"id":"c1","result":[{"id":"R17","title":"读书会"}]}<|end|>
<|assistant|><tool_call>{"id":"c2","name":"read","arguments":{"id":"R17"}}</tool_call>
<tool_call>{"id":"c3","name":"availability",
  "arguments":{"start":"2025-04-18T14:00:00+08:00","room":"Haitang-201"}}</tool_call><|end|>
<|tool|>{"id":"c3","result":{"available":true}}<|end|>
<|tool|>{"id":"c2","result":{"version":1,"room":"Haitang-101","note":"带笔记"}}<|end|>
\end{codeblock}

这个模板表达了两种多次调用。第一种是顺序调用：搜索返回编号以后，模型才能读取\code{R17}，于是调用与结果在助手消息和工具消息之间交替。第二种是\term{并行函数调用}（parallel function calling）：读取记录与查询时段互不依赖，模型在同一条助手消息中输出两段调用，各带一个调用编号；执行器可以同时执行，结果按完成先后返回，模型依靠编号而不是顺序把结果对应回调用。若没有编号，第二条工具消息先到时，模型就可能把时段查询结果当成记录内容。

运行时在解码过程中监视特殊词元：一旦出现\code{<tool\_call>}，就缓存后续词元直到\code{</tool\_call>}，再交给算法\ref{alg:nlp-dialogue-function-loop}那样的校验与执行流程，不合格的调用以错误消息返回给模型修正。也可以在调用片段内部启用约束解码，按JSON模式屏蔽非法词元（第\ref{sec:nlp-inference-grammar}小节），从源头保证片段可解析；参数值是否指向正确对象仍要另行检查（第\ref{sec:nlp-execution-tool-arguments}小节）。

有了模板，函数调用能力就可以用有监督微调获得：把多轮示范按模板序列化，在助手消息的词元上计算损失，包括思考文字、调用片段和最终回复；系统消息、用户消息和工具消息的词元则用式\eqref{eq:nlp-sft-masked-loss}的损失掩码排除。工具消息必须排除，因为它的内容由环境产生：若训练模型去预测工具结果，模型学到的恰是在没有真实返回时编造结果。示范从哪里来，是接下来两种方法要回答的问题。

\subsubsection{自监督工具调用学习：Toolformer}
\label{sec:nlp-execution-toolformer}

人工示范要回答两个难题：在一段文本的哪个位置调用工具才有帮助，以及调用写成什么参数。Toolformer的想法是让语言模型自己提出候选调用，执行后检查结果能否帮助预测后文；能帮助的调用保留下来，并入原文本，再用这些文本微调同一个模型\scite{nlp-execution-schick2023toolformer}。这样得到的训练数据只依赖少量调用示例和工具本身，不需要人工标注调用位置。

先约定调用的文本表示。一次调用$c=(a_c,i_c)$由工具名$a_c$和输入$i_c$组成，结果为$r$。用特殊标记把它线性化为
\[
 e(c)=\langle\mathrm{API}\rangle\,a_c(i_c)\,\langle/\mathrm{API}\rangle,\qquad
 e(c,r)=\langle\mathrm{API}\rangle\,a_c(i_c)\to r\,\langle/\mathrm{API}\rangle,
\]
前者只有调用，后者附上结果。原文使用的工具包括问答系统、计算器、维基百科搜索、机器翻译和日历，基座模型为GPT-J。

给定一段文本$\bm{x}=(x_1,\ldots,x_n)$和提示$P(\bm{x})$（提示中含该工具的若干调用示例，并要求模型在文本中插入调用），构造过程分四步。

第一步是采样调用位置。对每个位置$i$，计算模型在此处开始一次调用的概率
\[
 p_i=p_M\bigl(\langle\mathrm{API}\rangle\mid P(\bm{x}),x_{1:i-1}\bigr),
\]
只保留$p_i>\tau_s$的位置，若超过$k$个则取概率最高的$k$个。$\tau_s$是采样阈值。这一步用模型自己的判断筛掉大量不需要调用的位置。

第二步是采样调用。对每个保留位置$i$，以$[P(\bm{x}),x_1,\ldots,x_{i-1},\langle\mathrm{API}\rangle]$为前缀，最多采样$m$个调用$c_i^1,\ldots,c_i^m$，以$\langle/\mathrm{API}\rangle$为结束标记，没有生成结束标记的样本丢弃。

第三步是执行全部候选调用，得到结果$r_i$。工具必须能被批量自动执行，而且执行没有副作用。

第四步是按损失过滤。设$z$为放在文本之前的前缀，位置$i$之后的加权交叉熵为
\begin{equation}
 L_i(z)=-\sum_{j=i}^{n}w_{j-i}\log p_M(x_j\mid z,x_{1:j-1}),
 \label{eq:nlp-execution-toolformer-loss}
\end{equation}
其中权重$w_t$随距离$t$递减，原文取$\tilde w_t=\max(0,1-0.2t)$再归一化，只有紧随其后的五个词元计入。调用是否有用，比较三种前缀下的损失：
\begin{equation}
 L_i^+=L_i\bigl(e(c_i,r_i)\bigr),\qquad
 L_i^-=\min\Bigl(L_i(\varepsilon),\,L_i\bigl(e(c_i,\varepsilon)\bigr)\Bigr),\qquad
 L_i^--L_i^+\geq\tau_f ,
 \label{eq:nlp-execution-toolformer-filter}
\end{equation}
其中$\varepsilon$为空序列，$e(c_i,\varepsilon)$表示只写调用、不给结果。只有满足不等式的调用才保留，$\tau_f$为过滤阈值。$L_i^-$取两者中较小的一个，是为了把收益归于结果$r_i$：调用的输入本身可能已经透露了后文，例如输入中出现了要预测的词；若只与“不调用”比较，这种调用也会被误认为有用。与“只写调用”的损失也作比较，就要求结果在调用输入之外再提供信息。还有一处实现细节：调用与结果作为前缀放在文本之前，而不是插入位置$i$，因为此时模型尚未见过含调用的文本，插在中间会打断它熟悉的文本分布，使损失失去可比性。

过滤后，把保留的调用插回原位置得到$\bm{x}^*=(x_{1:i-1},e(c_i,r_i),x_{i:n})$，汇成增广语料，用普通的语言模型目标微调$M$。推断时，模型照常解码；一旦生成“$\to$”，运行时暂停解码，执行调用，把结果与$\langle/\mathrm{API}\rangle$写入序列，再继续生成。

下面用一个教学算例走一遍过滤准则，数值均为构造。文本为“读书会4月18日举行，距4月15日发布通知还有3天，请按时参加”，在“还有”之后的位置$i$采样到计算器调用$c$=\code{Calculator(18-15)}，结果$r=3$。五个权重为$\tilde w=(1,0.8,0.6,0.4,0.2)$，归一化后为$(1/3,4/15,1/5,2/15,1/15)$。设后续五个词元“3、天、，、请、按”在三种前缀下的负对数概率分别为
\[
\begin{aligned}
 \text{不调用：}&(2.40,0.30,0.20,1.10,0.90), & L_i(\varepsilon)&=1.127,\\
 \text{只写调用：}&(2.10,0.30,0.20,1.10,0.90), & L_i(e(c,\varepsilon))&=1.027,\\
 \text{调用加结果：}&(0.10,0.20,0.20,1.10,0.90), & L_i(e(c,r))&=0.333.
\end{aligned}
\]
以第一行为例，$2.40/3+0.30\times4/15+0.20/5+1.10\times2/15+0.90/15=0.800+0.080+0.040+0.147+0.060=1.127$。于是$L_i^-=1.027$，$L_i^--L_i^+=0.694$；取教学阈值$\tau_f=0.5$，这次调用保留。若在同一位置采样到问答调用，返回的是“海棠楼201”，与后文数字无关，设其带结果损失为$0.993$、只写调用损失为$1.060$，则$L_i^-=1.060$，差值仅$0.067$，被丢弃。

Toolformer与第\ref{sec:nlp-reasoning-rft}小节的拒绝采样微调同属“模型自己生成、按准则过滤、再微调”的自训练路线，区别在于过滤准则：拒绝采样按答案正误，Toolformer按语言模型损失的下降。增广语料除插入的调用外与预训练文本同源，原文借此让模型在学会调用的同时保留通用的语言建模能力。它的局限也来自这套构造：每个调用独立采样，模型学不到把一个工具的输出喂给另一个工具的链式调用；准则只看紧随其后的几个词元，只奖励能直接帮助预测后文的调用；工具必须在大规模文本上批量执行，因此只适用于没有副作用的查询类工具。本例的更新接口一旦执行就改动预约，无法用这种方式在预训练文本上试调用。

\subsubsection{工具使用指令数据的构造}
\label{sec:nlp-execution-tool-data}

Toolformer覆盖的是少量通用工具和单次调用。实际应用面对的是成千上万个接口，一条用户指令常要多次调用、组合多个接口。这需要大规模的“指令—调用轨迹”数据，人工编写不现实，于是改用强模型合成。

ToolLLM的数据构造分三步\scite{nlp-execution-qin2024toolllm}。第一步收集接口：从RapidAPI平台收集一万六千多个真实REST接口及其文档，按类别与集合组织。第二步生成指令：从接口集合中采样工具组合，分单工具、同类别多工具、同集合多工具三种情形，把采样接口的文档交给ChatGPT，要求它写出需要这些接口才能完成的指令，并列出相关接口；这些“指令—相关接口”对日后也用于训练接口检索器。第三步标注调用轨迹：让ChatGPT以函数调用方式真实调用接口，逐步完成指令。

第三步的困难在于，长轨迹中任何一步出错，ReAct式的单链生成就会整条失败，许多指令标注不出可用轨迹。ToolLLM为此提出\term{深度优先决策树搜索}（depth-first search-based decision tree, DFSDT）：把轨迹看成一棵树，每个节点是一个部分轨迹；在节点处由模型生成下一步调用作为子节点；若某条分支走到失败，或模型主动选择放弃，就回溯到父节点，生成一个与已试子节点不同的新子节点；找到第一条以有效答案结束的路径即停止。算法\ref{alg:nlp-execution-dfsdt}写出这一过程。

\begin{algorithm}[DFSDT调用轨迹标注]
\label{alg:nlp-execution-dfsdt}
\begin{algorithmic}[1]
\Require 指令$q$，接口集合$\mathcal{A}$及文档，标注模型$M$，调用预算$B$
\Ensure 一条以有效答案结束的轨迹，或标注失败
\Function{Expand}{$h$}\Comment{$h$为当前部分轨迹}
  \State $C\gets\emptyset$\Comment{本节点已试过的子节点}
  \While{$B>0$}
    \State 以$q,\mathcal{A},h$及已试子节点$C$为条件，令$M$生成与$C$不同的下一步$a$
    \If{$a$为给出答案}
      \State \Return $h\oplus a$
    \ElsIf{$a$为放弃}
      \State \Return 失败\Comment{回溯到父节点}
    \EndIf
    \State $o\gets$执行$a$；$B\gets B-1$；$C\gets C\cup\{a\}$
    \State $\tau\gets$\Call{Expand}{$h\oplus a\oplus o$}
    \If{$\tau$不是失败}
      \State \Return $\tau$
    \EndIf
  \EndWhile
  \State \Return 失败
\EndFunction
\State \Return \Call{Expand}{空轨迹}
\end{algorithmic}
\end{algorithm}

与ReAct式单链相比，DFSDT能从错误分支退回重试，提高了困难指令的标注成功率；与广度优先搜索相比，它只需要找到一条可行路径，不必展开同层所有分支，调用开销较小。生成新子节点时把已试子节点写进提示，是为了避免模型反复生成同一个失败调用。标注得到的轨迹用于微调LLaMA得到ToolLLaMA，评价时用模型评判的通过率与胜率衡量。

Gorilla面向的是机器学习模型库的接口调用，从TorchHub、TensorFlow Hub和HuggingFace收集接口文档，用自指令方法为每个接口生成指令\scite{nlp-execution-patil2024gorilla}。它的关键做法是\term{检索感知训练}（retriever-aware training）：训练时在指令后附上检索到的接口文档，写成“参考这份接口文档”的形式，让模型学会依据当前文档写调用，而不是只凭记忆。接口版本更新、参数改名后，测试时换入新文档即可，不必重新训练。评价时，Gorilla把生成的代码解析为抽象语法树，检查参考调用是否作为子树出现在其中；若生成的接口名不在接口库中，则记为幻觉。与运行代码相比，语法树匹配不需要执行环境，代价低，但只核对接口名与关键参数，不判断参数取值在语义上是否正确。

表\ref{tab:nlp-execution-tool-learning}对照三种获得调用能力的方式。合成数据的问题集中在两点：指令由文档反推，可能不像真实用户的需求；教师模型给出的轨迹可能本身有错，而按“得到答案”过滤也不保证答案正确。合成数据扩大了覆盖，但覆盖到的工具仍是文档描述清楚、能自动调用的那一部分。

\begin{table}[htbp]
\centering
\small
\begin{tabular}{p{17mm}p{24mm}p{28mm}p{26mm}}
\toprule
方法 & 监督来源 & 调用结构 & 过滤或评价 \\
\midrule
示范微调 & 人工或规则编写的多轮示范 & 任意，受示范覆盖限制 & 人工审核 \\
Toolformer & 模型在文本中自行插入调用 & 少量工具，单次、互不链接 & 后文损失下降$\geq\tau_f$ \\
ToolLLM & 强模型生成指令并搜索轨迹 & 大量真实接口，多步多工具 & DFSDT找到有效答案；模型评判 \\
Gorilla & 自指令生成指令，附检索文档 & 单次调用模型库接口 & 语法树子树匹配；幻觉率 \\
\bottomrule
\end{tabular}
\caption{获得函数调用能力的几种数据来源。各方法改变的是监督从哪里来、能覆盖怎样的调用结构，以及用什么准则筛选或评价数据。}
\label{tab:nlp-execution-tool-learning}
\end{table}

\subsection{工具检索与参数生成}
\label{sec:nlp-execution-tools}

模型具备调用能力之后，每一步仍要决定调用哪个工具、填什么参数。前者在工具很多时变成检索加过滤问题，后者要把每个值绑定到可靠来源；两者的错误都可能通过格式检查，需要按契约与对象身份另行核对。

\subsubsection{工具检索与动作选择}
\label{sec:nlp-execution-tool-selection}

教学目录有五个工具：只读的\code{search}支持按名称或按活动编号与用户查找，\code{read}读取单条记录，\code{availability}检查时段，\code{update}只改时间和房间且不发通知，\code{notify}向参与者发送消息。检索的查询应当是“当前还缺什么信息或变化”，而不是整段用户要求。尚无对象编号时应选搜索；已确认编号但不知道版本时应选读取；目标与前置条件都已确认，更新工具才成为候选。\code{notify}在语义上与“预约修改”高度相关，却违反本例的否定约束。

工具多时，常用描述嵌入做第一轮召回：以编码器$E$把查询$q$和工具描述$d_k$映射为向量，按余弦相似度$s(q,d_k)=\cos\bigl(E(q),E(d_k)\bigr)$取前$K$个候选写入上下文，再由模型在候选中选择。ToolLLM的接口检索器就用“指令—相关接口”对训练句向量编码器，以相关接口为正例、抽取的若干其他接口为负例做对比学习，目标与式\eqref{eq:nlp-matching-infonce}同类。召回之后还须过滤：调用者没有权限、缺少必需输入或副作用与约束冲突的工具，即使相似度最高也不能执行。模型生成目录外的工具名时，运行时应报告不可用，而不是拼凑一个接口。

名称相似的工具不能互换，需要比较的是契约。若真实系统的更新接口强制给所有参与者发消息，本例就没有满足约束的动作，正确做法是说明冲突或请求用户重新授权，而不是选一个“最接近”的工具。

\subsubsection{函数参数生成与模式约束}
\label{sec:nlp-execution-tool-arguments}

动作确定以后，参数要从有来源的字段绑定，而不是只把JSON填满。搜索结果提供\code{id}，刚读取的记录提供\code{expected\_version}，通知提供目标时间和房间，控制器生成任务内唯一的请求号。一个正常的更新请求为：

\begin{codeblock}[label={code:nlp-execution-tool-request}]{本地模拟工具的结构化参数}
{
  "id": "R17",
  "expected_version": 1,
  "start": "2025-04-18T14:00:00+08:00",
  "room": "Haitang-201",
  "request_id": "task7-goal1-v1"
}
\end{codeblock}

按代码\ref{code:nlp-execution-update-schema}的模式做约束解码，可以消除缺字段、多字段和类型错误，但对象绑定、时区与日历仍要逐项检查：解析带时区的时间，查询编号与版本，并把记录所属用户与活动编号同当前目标比对。本地接口核验中，\code{R404}得到不存在错误，字符串版本\code{"1"}得到模式错误，2月30日得到日期错误，过期版本得到冲突。更隐蔽的是两组合法候选：把日期填为通知发布日期4月15日，或把编号填为同名的\code{R18}。这些值都能通过模式与日期解析，若权限允许，接口会成功修改错误的目标。因此调用器在写前核对目标身份，结束判定再读取正确对象；参数文本与接口回执分别保存。

\subsection{计划与行动选择}
\label{sec:nlp-execution-planning}

多步任务要决定行动的顺序，以及观察出乎意料时怎样调整。预算至少包括模型调用次数、工具调用次数和总时间，失败调用与回查同样计入。三种控制结构在这些预算上各有取舍（图\ref{fig:nlp-execution-planning-patterns}）：先计划后执行只需很少的模型调用，计划也便于事先检查，但难以应对意外的观察；ReAct每一步都依据最新观察决定下一步，灵活，但模型调用多、上下文不断增长；依赖图让互不依赖的子任务并行，缩短等待时间。

\begin{figure*}[htbp]
  \centering
  % 第九卷 14.3.3：三种计划与行动组织方式（跨栏，约160 mm宽）。
% FigureModel：语言模型调用；FigureOutput：工具执行与调度；虚线：变量引用或重新规划。
\begin{tikzpicture}[
  x=1mm,y=1mm,
  figure picture,
  model/.style={figure node,draw=FigureModel,fill=FigureModelFill,
    minimum width=38mm,align=center,minimum height=11mm,inner xsep=2pt,line width=.8pt},
  tool/.style={figure node,draw=FigureOutput,fill=FigureOutputFill,
    minimum height=8mm,inner xsep=2pt},
  title/.style={font=\figurefont\small,text=FigureInk,align=center},
  ref/.style={figure flow,dash pattern=on 2pt off 1.5pt}
]
  % (a) ReAct
  \node[title] at (22,64) {（a）ReAct：逐步交替};
  \node[model] (ra) at (22,48) {语言模型\\Thought$_t$与Action$_t$};
  \node[tool,minimum width=30mm] (rt) at (22,22) {工具执行};
  \draw[figure flow] ([xshift=-7mm]ra.south) --
    node[figure edge label,left=1pt] {Action$_t$} ([xshift=-7mm]rt.north);
  \draw[figure flow] ([xshift=7mm]rt.north) --
    node[figure edge label,right=1pt] {Observation$_t$} ([xshift=7mm]ra.south);
  \node[figure note] at (22,7) {$T$次工具调用需$T{+}1$次模型调用\\每次重读提示与全部历史};

  % (b) ReWOO
  \node[title] at (74,64) {（b）ReWOO：先规划后执行};
  \node[model] (wp) at (74,48) {规划器\\一次写出含\#E$_k$的计划};
  \node[tool,minimum width=11mm] (w1) at (60,31) {\#E$_1$};
  \node[tool,minimum width=11mm] (w2) at (74,31) {\#E$_2$};
  \node[tool,minimum width=11mm] (w3) at (88,31) {\#E$_3$};
  \node[model] (ws) at (74,14) {求解器\\读计划与证据作答};
  \foreach \w in {w1,w2,w3} {
    \draw[figure flow] (wp.south -| \w) -- (\w.north);
    \draw[figure flow] (\w.south) -- (ws.north -| \w);
  }
  \draw[ref] (w1.east) -- (w2.west);
  \node[figure note] at (74,3) {工作器按变量引用代入结果\\共2次模型调用};

  % (c) LLMCompiler
  \node[title] at (126,64) {（c）LLMCompiler：依赖图并行};
  \node[model] (cp) at (126,48) {函数调用规划器\\输出任务依赖图};
  \node[tool,minimum width=34mm,align=center] (cf) at (126,34) {任务获取单元\\依赖满足即派发};
  \node[tool,minimum width=15mm] (c1) at (117,21) {执行器};
  \node[tool,minimum width=15mm] (c2) at (135,21) {执行器};
  \node[model] (cj) at (126,7) {汇合\\作答或重新规划};
  \draw[figure flow] (cp.south) -- (cf.north);
  \draw[figure flow] (cf.south -| c1) -- (c1.north);
  \draw[figure flow] (cf.south -| c2) -- (c2.north);
  \draw[figure flow] (c1.south) -- (cj.north -| c1);
  \draw[figure flow] (c2.south) -- (cj.north -| c2);
  \draw[ref] (cj.east) -- (148,7) -- node[figure edge label,right=1pt,text width=8mm] {重新\\规划} (148,48) -- (cp.east);
\end{tikzpicture}

  \caption{三种计划与行动的组织方式。（a）ReAct中模型每写出一个行动就等待真实观察，再决定下一步；（b）ReWOO由规划器一次写出含证据变量的计划，工作器依次执行并代入变量（虚线表示$\#E_2$引用$\#E_1$），求解器最后作答；（c）LLMCompiler的规划器输出任务依赖图，任务获取单元把依赖已满足的任务派给执行器并行执行，汇合步骤决定作答或重新规划。}
  \label{fig:nlp-execution-planning-patterns}
\end{figure*}

\subsubsection{计划与执行的分离}
\label{sec:nlp-execution-plan-execute}

零样本思维链（第\ref{sec:nlp-execution-cot}小节）常见的一类错误是漏掉步骤：模型直接写推理，写到一半忘了题目中的某个条件。Plan-and-Solve提示把一次推理拆成两段，先让模型理解问题并制订分步计划，再按计划逐步执行\scite{nlp-execution-wang2023planandsolve}；其增强版本还要求模型先抽取相关变量及其数值、注意中间结果的计算。它仍是一次生成，计划与执行写在同一段文字中，作用是让模型在执行前显式列出子目标，减少遗漏。

把这一思想用于工具调用，就得到ReWOO\scite{nlp-execution-xu2023rewoo}。它把系统拆成三个角色：\term{规划器}（planner）一次写出完整计划，每一步是一句说明加一个工具调用，调用结果存入证据变量$\#E_k$，后续步骤的输入可以引用前面的变量；\term{工作器}（worker）按顺序执行各步调用，把变量引用替换为实际结果；\term{求解器}（solver）读入任务、计划与全部证据，写出最终答案。规划时模型看不到任何观察，这正是方法名称“无观察推理”的含义。代码\ref{code:nlp-execution-rewoo-plan}是本例按此格式写出的计划。

\begin{codeblock}[label={code:nlp-execution-rewoo-plan}]{ReWOO格式的计划（教学示例）}
Plan: 按用户与活动中文名称查找预约。
#E1 = search[user="U7", name="机器学习读书会"]
Plan: 读取所找到记录的当前版本与字段。
#E2 = read[id=#E1.id]
Plan: 查询目标时段与房间是否可用。
#E3 = availability[start="2025-04-18T14:00:00+08:00", room="Haitang-201"]
Plan: 时段可用时按读取的版本更新时间和房间，不发通知。
#E4 = update[id=#E1.id, expected_version=#E2.version, ...]
\end{codeblock}

分离的收益首先体现在上下文开销上。设提示（任务说明、工具描述与示例）长$P$个词元，每一步的调用与观察平均增加$h$个词元，共$T$步。ReAct式交替要调用模型$T+1$次，第$t$次读入$P+(t-1)h$个词元；ReWOO只调用规划器和求解器各一次，求解器读入提示与全部证据。两者读入的词元总数分别约为
\begin{equation}
 C_{\mathrm{ReAct}}=(T+1)P+h\,\frac{T(T+1)}{2},\qquad
 C_{\mathrm{ReWOO}}=P_{\mathrm{plan}}+P_{\mathrm{solve}}+Th .
 \label{eq:nlp-execution-token-cost}
\end{equation}
前者随步数平方增长，后者线性增长。教学数值取$P=P_{\mathrm{plan}}=P_{\mathrm{solve}}=1500$、$h=200$、$T=5$，得$C_{\mathrm{ReAct}}=6\times1500+200\times15=12000$，$C_{\mathrm{ReWOO}}=1500+1500+1000=4000$。由于规划器不必读入观察，原文还把规划能力蒸馏到较小的模型上。

代价是计划无法根据观察调整。模拟目录存的是英文名称\code{ML reading group}，第一步按中文名称查询返回空列表，$\#E1$没有\code{id}字段，引用它的第二步和第四步都无法执行。纯粹的ReWOO只能把错误交给求解器，由它说明未完成。本例做一处扩展：工作器发现变量解析失败时，只把产生$\#E1$的子目标退回规划器，改按活动编号\code{ML-0418}查找，不依赖$\#E1$的$\#E3$及其结果保留。计划中含写操作时，写入前还须核对前面的证据确实满足条件，否则一次性计划会在版本已变或时段已占时照常提交。

\subsubsection{判断、行动与观察交替：ReAct}
\label{sec:nlp-execution-react}

先规划的方法在观察出乎意料时缺乏调整手段；只写思维链的方法没有外部信息，容易编造事实；只输出行动、不写推理的方法又难以分解目标、跟踪进度。ReAct把两者交织起来：在行动空间$\mathcal{A}$之外加入语言空间$\mathcal{L}$，模型在每一步可以写一段思考（thought），也可以发出一个行动\scite{yao2023react}。思考不改变环境，只改变上下文，用来分解目标、记录进度、解释观察和决定下一步。形式上，设$g$为目标，历史$h_t=(\mathit{th}_1,a_1,o_1,\ldots,\mathit{th}_{t-1},a_{t-1},o_{t-1})$，模型在每一步生成
\[
 (\mathit{th}_t,a_t)\sim\pi_{\bm{\theta}}(\,\cdot\mid g,h_t),\qquad a_t\in\mathcal{A},
\]
执行器执行$a_t$得到观察$o_t$。原文用少样本示范的轨迹作提示，示范逐行写成“Thought $t$”“Action $t$”“Observation $t$”三类行；在维基百科问答任务上，行动空间是\code{Search[实体]}、\code{Lookup[字符串]}和\code{Finish[答案]}。模型写完一行Action后，运行时以“Observation”为停止序列中断生成，执行行动，把真实返回写成Observation行，再让模型继续。Observation行永远由执行器写入，模型自己写出的“观察”不算数。代码\ref{code:nlp-execution-react-trace}是本例按这一格式记录的轨迹。

\begin{codeblock}[label={code:nlp-execution-react-trace}]{ReAct格式的执行轨迹（教学模拟）}
Question: 把U7在活动ML-0418的预约改到通知中的时间和地点，不要通知，保留备注。
Thought 1: 先要找到预约编号。按活动中文名称搜索。
Action 1: search[user="U7", name="机器学习读书会"]
Observation 1: []
Thought 2: 中文名称没有命中，改用确认过的活动编号。
Action 2: search[user="U7", event="ML-0418"]
Observation 2: [{"id": "R17", "event": "ML-0418"}]
Thought 3: 找到R17，读取当前版本与字段。
Action 3: read[id="R17"]
Observation 3: {"version": 1, "start": "2025-04-17T14:00:00+08:00",
               "room": "Haitang-101", "note": "带笔记"}
Thought 4: 修改前确认目标时段可用。
Action 4: availability[start="2025-04-18T14:00:00+08:00", room="Haitang-201"]
Observation 4: {"available": true}
Thought 5: 条件满足，按版本1更新时间和房间，不发通知。
Action 5: update[id="R17", expected_version=1, ...]
Observation 5: {"status": "committed", "new_version": 2}
Thought 6: 回执只说明请求已提交，还要读回字段。
Action 6: read[id="R17"]
Observation 6: {"version": 2, "start": "2025-04-18T14:00:00+08:00",
               "room": "Haitang-201", "note": "带笔记"}
Thought 7: 时间、地点已改，备注未变，任务完成。
Action 7: finish[R17已改到4月18日14:00海棠楼201，未发送通知，备注保留]
\end{codeblock}

这条轨迹共有六次工具调用。第1次的空结果使思考转向按编号查找，这正是ReAct相对一次性计划的优势：调整只发生在出问题的那一步。若把空结果误写成“已找到”，后续文字仍会很流畅，却已失去对象证据；若调用上限为4，执行器必须停在时段查询之后，报告尚未修改。配套程序中的动作由确定性规则选择，轨迹只展示格式与控制流，不能据此估计ReAct模型的成功率。

算法\ref{alg:nlp-execution-react}在这种交替上补充本章的预算与异常处理，这些工程条件不属于原论文设计。完成请求要先转成一次读取核验；写入结果未知时，下一步必须是查询原请求或读取对象，不允许另一个写动作越过它。

\begin{algorithm}[带预算与结果核验的行动循环]
\label{alg:nlp-execution-react}
\begin{algorithmic}[1]
\Require 目标$g$，工具目录，初始状态$S$，调用预算$B$，决策与时间上限
\Ensure 已核验完成，或附当前证据的未完成状态
\While{决策和时间预算未耗尽}
  \State 接收用户更正，更新目标版本，标记受影响的计划与事实
  \If{$B=0$}
    \State \Return 预算耗尽，附未核验项与结果未知的操作
  \EndIf
  \If{存在结果未知的写入}
    \State $a\gets$查询原请求状态或读取实际对象
  \Else
    \State 依据$g,S$写出思考并选择$a$；完成请求改为结束核验
  \EndIf
  \State 检查参数、身份、权限与前置条件；不符则有限次修正或退出
  \State $o\gets$实际调用$a$；$B\gets B-1$；保存参数、返回与时间
  \State 按事件更新$S$；写入超时记为未知，不记为成功或失败
  \If{最新核验观察满足$g$及保持条件，且无未知操作}
    \State \Return 已核验完成及其证据
  \EndIf
\EndWhile
\State \Return 达到上限，附已发生的变化与剩余条件
\end{algorithmic}
\end{algorithm}

ReAct与前后方法的关系可以这样概括：它是思维链加上外部行动，也是第\ref{sec:nlp-execution-swe-agent}小节SWE-agent在仓库修复中所用的基本循环。原文除提示外，还用模型自举得到的成功轨迹微调较小的模型，这条路线在第\ref{sec:nlp-execution-trajectory-sft}小节展开。它的代价已由式\eqref{eq:nlp-execution-token-cost}给出：模型调用次数与工具调用次数同阶，上下文随步数增长；模型还可能陷入重复同一行动的循环，或写出与随后行动不一致的思考，因此需要调用上限和重复检测。

\subsubsection{子目标依赖图与并行调用}
\label{sec:nlp-execution-subgoals}

ReAct每一步只发出一个行动，ReWOO的工作器按计划顺序执行，二者都把互不依赖的调用串行化了。工具调用的耗时常由网络与外部服务决定，串行等待会累积成主要延迟。LLMCompiler借用编译器的思路，把计划写成任务依赖图，再按依赖关系并行执行\scite{nlp-execution-kim2024llmcompiler}。它由三个部分组成：\term{函数调用规划器}（function calling planner）输出任务列表，每个任务是一次工具调用，输入可以用占位符\code{\$k}引用第$k$个任务的输出，占位符关系即依赖边；\term{任务获取单元}（task fetching unit）维护各任务状态，一旦某任务的全部前驱完成，就把占位符替换为实际输出并派发；\term{执行器}（executor）并行执行派来的调用。全部任务完成后，由模型汇合结果，决定直接作答还是重新规划；规划器还可以流式输出任务，前几个任务在计划写完前就开始执行。

设任务依赖图为$\mathcal{G}=(V,E)$，任务$v$的执行时间为$t_v$。顺序执行的总时间是$\sum_{v\in V}t_v$；按依赖并行时，总时间由图中耗时最长的路径决定，即\term{关键路径}（critical path）：
\begin{equation}
 T_{\mathrm{seq}}=\sum_{v\in V}t_v,\qquad
 T_{\mathrm{par}}=\max_{\text{路径}\,\rho\subseteq\mathcal{G}}\ \sum_{v\in\rho}t_v .
 \label{eq:nlp-execution-critical-path}
\end{equation}
并行缩短的是等待时间，调用次数不变。算法\ref{alg:nlp-execution-task-fetch}写出任务获取单元的调度。

\begin{algorithm}[任务获取单元的依赖调度]
\label{alg:nlp-execution-task-fetch}
\begin{algorithmic}[1]
\Require 任务依赖图$\mathcal{G}=(V,E)$，各任务的调用模板
\Ensure 各任务的输出，或失败任务及其后继
\State $R\gets\{v\in V:v\text{没有前驱}\}$；$D\gets\emptyset$\Comment{$R$为就绪集合，$D$为已完成集合}
\While{存在未完成任务}
  \For{$v\in R$}
    \State 把$v$输入中的占位符替换为前驱的实际输出，异步派发给执行器
  \EndFor
  \State $R\gets\emptyset$；等待任一任务$u$返回
  \If{$u$失败}
    \State 标记$u$在$\mathcal{G}$中的全部后继为不可执行，交汇合步骤处理
  \Else
    \State $D\gets D\cup\{u\}$
    \State $R\gets R\cup\{v:(u,v)\in E,\ v\text{的前驱全在}D\text{中}\}$
  \EndIf
\EndWhile
\end{algorithmic}
\end{algorithm}

为展示并行，另构造一个会前准备预约\code{R21}：用户明确允许4月18日13:00使用海棠楼202或203，优先202，读书会本身仍在201。任务图为

\begin{codeblock}[label={code:nlp-execution-compiler-plan}]{LLMCompiler格式的任务列表（教学示例）}
1. read(id="R21")
2. availability(room="Haitang-202", start="2025-04-18T13:00:00+08:00")
3. availability(room="Haitang-203", start="2025-04-18T13:00:00+08:00")
4. choose_room(prefer=["Haitang-202", "Haitang-203"], results=[$2, $3])
5. update(id="R21", expected_version=$1.version, room=$4, ...)
6. join()
\end{codeblock}

任务1、2、3没有前驱，可以同时派发；任务4依赖2与3，任务5依赖1与4。设读取耗时0.5秒，每次时段查询0.8秒，选择为本地计算不计时，更新0.6秒。顺序执行需$0.5+0.8+0.8+0.6=2.7$秒；并行时关键路径为“查询—选择—更新”，需$0.8+0.6=1.4$秒，两者都是四次服务调用。反过来，若202一查就可用，按优先规则的顺序策略可以不查203，只用三次调用；预先并行则多付一次调用，换取的是不必等待。并行是否划算，取决于延迟与调用成本哪个更紧。

依赖图还使更正的影响可以计算。用户把日期改为19日后，依赖日期的任务2、3及其后继4、5都要失效，按图的可达关系即可找出；任务1读取的记录与日期无关，结果保留。若只有任务2故障，也只需重做任务2及其后继。规划器的责任因此更重：漏写一条依赖边，例如让更新不依赖时段查询，执行器就可能在查询返回前提交写入。有副作用的任务不宜投机地提前并行，只读调用才适合积极并行。表\ref{tab:nlp-execution-planning-compare}汇总三种控制结构。

\begin{table}[htbp]
\centering
\small
\begin{tabular}{p{21mm}p{21mm}p{25mm}p{28mm}}
\toprule
方法 & 模型调用次数 & 对意外观察的响应 & 主要代价 \\
\midrule
ReWOO & 固定2次 & 计划内无法响应，交给求解器或重规划 & 变量解析失败使后续步骤作废 \\
ReAct & $T{+}1$次 & 每步依据最新观察调整 & 读入词元随步数平方增长，串行等待 \\
LLMCompiler & 规划与汇合各1次，可重规划 & 汇合时整体判断，可重规划 & 依赖声明错误导致竞争；可能多付调用 \\
\bottomrule
\end{tabular}
\caption{三种计划与行动结构的对照。$T$为工具调用步数；模型调用次数不含失败重试。}
\label{tab:nlp-execution-planning-compare}
\end{table}

\subsection{观察、状态与任务记忆}
\label{sec:nlp-execution-state-memory}

轨迹变长以后，执行器要有选择地保留信息：当前状态只保存仍影响任务的事实，上下文装不下的历史要压缩或按需取回，跨回合尝试还要带上失败经验。三个小节分别处理这三类信息。

\subsubsection{事件更新与结构化任务状态}
\label{sec:nlp-execution-state-update}

状态中的每个字段都应能回答“谁说的，针对哪个对象，在哪个版本成立”。把状态组织为$S=(g^{(v)},I,F,H,O)$：$g^{(v)}$是第$v$版用户目标，$I$为已确认的对象身份，$F$为带来源和版本的观察事实，$H$为尚未确认的假设，$O$为操作及其状态。每个事件$e_t$携带来源、对象编号、观察时间以及环境版本或请求号；更新$S_{t+1}=U(S_t,e_t)$按接口规则进行，不需要训练模型去猜写入是否发生。

正常轨迹中，第3次调用后$F$保存\code{R17}版本1的旧字段，$O$为“未提交”；第5次得到回执后，$O$记为“已提交，回执版本2”，但$F$中的字段仍待回读；第6次读取后$F$才有新快照。写入超时时，$O$记为“结果未知”，期望字段不复制进$F$。若其他入口在第5次与第6次之间改了备注，版本3的读取覆盖字段快照，版本2的回执留作审计；反过来，迟到的版本2响应不能覆盖已知的版本3。用户说“改到19日”更新的是$g$，不意味着环境已处于19日，两者分属不同字段。

\subsubsection{历史摘要与检索式记忆}
\label{sec:nlp-execution-history}

完整轨迹直接放进上下文能保留细节，却占用位置并混入过期信息。压缩历史有两条路：写成摘要，或存入外部记忆、按需检索。摘要至少要保存目标版本、对象编号、否定条件、已确认字段、未知操作及原始结果的定位。若摘要只写“预约改到18日”，就丢掉了“不发通知”和“保留备注”，执行器可能选择一个会通知全员的接口。较可靠的写法是“用户U7，活动ML-0418，目标18日14:00与201，不通知；R17版本1已读取；写入尚未提交”，并附各字段的来源事件号。遇到不确定的内容，摘要应保留“待核对”，不能在压缩时把“可能是R17”写成“已确认R17”。

检索式记忆要回答的问题是：记忆库很大时，当前该取回哪几条。Generative Agents把智能体经历的事件逐条写成自然语言记录，组成\term{记忆流}（memory stream），每条记录带创建时间与最近访问时间\scite{nlp-execution-park2023generative}。取回时，它对每条记忆$m$按三个分量打分：新近度按距上次访问的时间指数衰减，重要性由模型在1到10之间打分（日常琐事低、关键事件高），相关性是记忆与当前查询$q$的嵌入余弦相似度。三个分量各自在候选集合上做最小—最大归一化到$[0,1]$，再加权求和：
\begin{equation}
 \mathrm{score}(m,q)=\alpha_{\mathrm{rec}}\,\widehat{\gamma^{\Delta(m)}}
 +\alpha_{\mathrm{imp}}\,\widehat{\mathrm{imp}(m)}
 +\alpha_{\mathrm{rel}}\,\widehat{\cos\bigl(\bm{e}_m,\bm{e}_q\bigr)},
 \label{eq:nlp-execution-memory-score}
\end{equation}
其中$\Delta(m)$是距上次访问的时间，$\gamma<1$是衰减因子，帽子表示归一化，$\bm{e}_m,\bm{e}_q$为嵌入向量。原文以游戏内小时计时，取$\gamma=0.995$，三项权重均为1，按得分从高到低取能放进上下文的记忆。

用一个教学算例看三项如何相互作用。设当前查询为“更新R17的时间”，有三条候选记忆：$m_1$为“用户要求不要发送通知”，2小时前访问，重要性8，相关性0.62；$m_2$为“R18上次更新成功”，1小时前，重要性4，相关性0.71；$m_3$为“用户偏好下午活动”，30小时前，重要性3，相关性0.40。新近度为$0.995^2=0.990$、$0.995^1=0.995$、$0.995^{30}=0.860$，归一化后为$0.963$、$1$、$0$；重要性归一化后为$1$、$0.2$、$0$；相关性归一化后为$0.710$、$1$、$0$。总分分别为$2.673$、$2.200$、$0$。只按相关性，排第一的是另一个对象的$m_2$；加入重要性后，否定约束$m_1$排到第一。但$m_2$仍排第二，会被一起取回：打分只衡量文本是否值得读，不判断它是否适用于当前对象，取回后仍要按对象编号、时间与有效性过滤。原文还在重要性累计超过阈值时，让模型从近期记忆中归纳出更高层的\emph{反思}记录，也存入记忆流参与检索。

Generative Agents的检索规则是固定的：何时取、取多少由外部程序决定。MemGPT借用操作系统虚拟内存的思路，让模型自己管理记忆\scite{nlp-execution-packer2023memgpt}。它把上下文分为主上下文与外部上下文两层：主上下文相当于内存，含系统指令、一块可读写的工作上下文和一个先进先出的消息队列；外部上下文相当于磁盘，含保存全部历史消息的回忆存储和存放任意文本的归档存储。模型通过函数调用在两层之间搬运信息，例如把重要事实写入工作上下文、在归档存储中检索；队列接近上下文容量时，系统先向模型发出内存压力提示，让它保存要紧的信息，随后把较早的消息移出队列，并用递归摘要保留其要点。这种\term{虚拟上下文管理}（virtual context management）让取回决策随任务变化，代价是记忆操作本身也要消耗调用与上下文，模型漏存或误删时同样会丢信息。

两种机制都沿用第\ref{chap:nlp-classification-matching-organization}章的检索方法负责召回。评价记忆时，应同时检查关键条件的保留率、取回记录的身份正确率和最终执行结果；检索命中一条相似历史，只说明找到了文本，不说明当时的动作现在可以重用。

\subsubsection{反馈反思与回合记忆：Reflexion}
\label{sec:nlp-execution-reflexion}

一次尝试失败后，下一次尝试若只拿到原始轨迹，行动者要自己从长轨迹中找出关键错误。Reflexion让模型把失败写成一段简短的文字教训，作为下一回合的上下文\scite{nlp-execution-shinn2023reflexion}。它区分三个职责：\term{行动者}（actor）在环境中生成轨迹；评价器（evaluator）给出成败信号，可以是规则、单元测试、环境奖励或模型判断；\term{自我反思}（self-reflection）模型依据轨迹与评价写出反思。设第$j$次尝试的轨迹为$\tau_j$，评价为$r_j$，反思为
\[
 m_j=f_{\psi}(\tau_j,r_j),\qquad M_{j+1}=\mathrm{keep}(M_j\cup\{m_j\}),
\]
其中$\mathrm{keep}$只保留最近的少数几条。下一次尝试的行动者以原目标、新观察和$M_{j+1}$为条件。原文称之为“语言强化”：评价信号起到奖励的作用，但更新的是上下文中的记忆，而不是模型参数。评价通过或尝试次数用完时退出。

构造一条旧记忆：上一轮直接复用了历史编号\code{R404}，接口返回不存在。反思可以写成“上轮复用了历史编号，未核对当前活动；下一轮先按用户与活动编号重新查找，读取版本后再修改”，并链接到失败调用。下一轮只有真的先搜索、得到\code{R17}并完成核验，才能说反思起了作用。与另两种记忆条件相比：不保存记忆，容易再次复用旧编号；只保存原轨迹，要行动者自己定位错误；保存反思，提供了压缩后的修正规则，但反思也可能归因错误，例如把失败归于权限不足而建议换账号，因此外部错误码与反思器对原因的猜测应分别标注。

\subsection{智能体的训练}
\label{sec:nlp-execution-agent-training}

以上方法都靠提示驱动：基座模型能否写出合格的思考与调用，决定了智能体的上限。较小的模型按少样本ReAct提示工作时，常出现格式错误、重复行动和过早结束，长提示本身也占用上下文。训练可以直接提升多步工具使用。轨迹微调用成功轨迹做有监督学习，简单稳定，但只能模仿示范；多轮强化学习以任务结果为奖励在环境中探索，能超出示范，代价是奖励稀疏、轨迹长、环境交互昂贵。

\subsubsection{交互轨迹微调}
\label{sec:nlp-execution-trajectory-sft}

轨迹微调的数据流程是：用强模型在环境中按ReAct格式执行任务，记录完整轨迹，按任务结果过滤，再用保留的轨迹微调较小的模型。FireAct用GPT-4在问答任务上调用搜索接口生成轨迹，微调Llama等模型，并发现混合ReAct、思维链和Reflexion等不同格式的轨迹，能让模型按任务难度灵活选择策略\scite{nlp-execution-chen2023fireact}。AgentTuning从操作系统、数据库、知识图谱、网页购物等多个环境收集GPT-4轨迹，只保留奖励达标的轨迹组成智能体指令数据\scite{nlp-execution-zeng2024agenttuning}。

设一条轨迹为$\tau=(g,\mathit{th}_1,a_1,o_1,\ldots,\mathit{th}_T,a_T)$，序列化后为词元序列$(z_1,\ldots,z_{|\tau|})$。损失只计在模型应当生成的思考与行动词元上，目标、系统提示和观察词元用掩码$m_k$排除，与式\eqref{eq:nlp-sft-masked-loss}相同：
\begin{equation}
 L_{\mathrm{agent}}(\bm{\theta})=-\frac{1}{\sum_k m_k}\sum_{k=1}^{|\tau|}m_k\log p_{\bm{\theta}}(z_k\mid z_{<k}),\qquad
 m_k=\begin{cases}1,& z_k\in\mathit{th}_t\cup a_t,\\0,&\text{其他}.\end{cases}
 \label{eq:nlp-execution-trajectory-sft}
\end{equation}
观察词元仍留在序列中作为条件，只是不计损失；理由与工具消息相同，模型不应学习预测环境的返回。只用智能体轨迹微调，模型的通用能力会下降，AgentTuning因此采用\term{混合指令微调}（hybrid instruction tuning），以比例$\eta$混合智能体数据与通用对话数据：
\begin{equation}
 L(\bm{\theta})=\eta\,L_{\mathrm{agent}}(\bm{\theta})+(1-\eta)\,L_{\mathrm{general}}(\bm{\theta}),
 \label{eq:nlp-execution-hybrid-tuning}
\end{equation}
原文在7B模型上以0.1为步长搜索$\eta$，$\eta=0.2$在未见任务上效果最好。

按成功过滤再微调，与第\ref{sec:nlp-reasoning-rft}小节的拒绝采样微调同构，只是“答案正确”换成了“任务成功”；从强化学习的角度看，它是对成功示范的行为克隆。行为克隆的弱点也随之而来：训练只见过教师走过的成功状态，学生一旦偏离，例如第1次搜索返回了教师轨迹中没有出现的空结果，就进入训练数据没有覆盖的状态。过滤只保留成功轨迹，失败中的恢复动作恰好被丢掉了。这些问题促使研究者在环境中直接优化任务结果。

\subsubsection{面向多轮交互的强化学习}
\label{sec:nlp-execution-agent-rl}

把整条交互看成一个回合：状态是目标与截至当前的历史，行动是模型在一轮中生成的思考与调用，环境执行后返回观察，回合结束时由判定器给出奖励$R(\tau)$，例如任务成功记1、否则记0。目标是最大化期望奖励$J(\bm{\theta})=\mathbb{E}_{\tau\sim p_{\bm{\theta}}}[R(\tau)]$。这里用到策略梯度定理：期望奖励的梯度等于奖励加权的轨迹对数概率梯度的期望，$\nabla J=\mathbb{E}[R(\tau)\nabla\log p_{\bm{\theta}}(\tau)]$，减去不依赖行动的基线不改变期望，只降低方差；一般推导与PPO等算法可参考\BookChapterRef{chap:rl-learning}{卷七“未知模型的求解”章}。

多轮交互的特殊之处在于轨迹概率由两类因子组成：
\begin{equation}
 p_{\bm{\theta}}(\tau)=\prod_{t=1}^{T}\pi_{\bm{\theta}}(\mathit{th}_t,a_t\mid g,h_t)\;P_{\mathrm{env}}(o_t\mid g,h_t,a_t).
 \label{eq:nlp-execution-trajectory-prob}
\end{equation}
环境转移$P_{\mathrm{env}}$不含参数，取对数求梯度时整项消失，$\nabla\log p_{\bm{\theta}}(\tau)=\sum_t\nabla\log\pi_{\bm{\theta}}(\mathit{th}_t,a_t\mid g,h_t)$。展开到词元层面，只有模型生成的词元贡献梯度，观察词元不计入。这正是实现中“观察词元掩码”的来源：它不是经验技巧，而是策略梯度在含环境步的轨迹上的直接推论。若不加掩码，等于把环境返回当成模型的行动去强化，梯度方向就错了。

第\ref{sec:nlp-reasoning-grpo}小节的GRPO可以直接推广到多轮。对同一任务从相同初态采样$G$条轨迹$\tau_1,\ldots,\tau_G$，得到奖励$r_1,\ldots,r_G$，按式\eqref{eq:nlp-reasoning-group-advantage}计算组内标准化优势$\widehat A_i$，赋给第$i$条轨迹中所有模型生成的词元。与单轮目标式\eqref{eq:nlp-reasoning-grpo-objective}相比，唯一的改动是把对回答全部词元的平均换成只对模型生成词元的掩码平均：记$\ell_{i,k}(\bm{\theta})$为式\eqref{eq:nlp-reasoning-grpo-objective}中第$i$条轨迹第$k$个词元上的裁剪项与KL项之和，$m_{i,k}$为该词元的掩码，则
\begin{equation}
 J(\bm{\theta})=\mathbb{E}\Biggl[\frac{1}{G}\sum_{i=1}^{G}\frac{1}{\sum_k m_{i,k}}\sum_{k}m_{i,k}\,\ell_{i,k}(\bm{\theta})\Biggr].
 \label{eq:nlp-execution-multiturn-grpo}
\end{equation}
观察词元仍作为条件参与计算概率比，只是不进入求和。Search-R1是这一做法的实例：模型在推理中发出搜索请求，检索结果插回序列，训练时对检索结果词元做掩码，以最终答案是否与参考精确匹配为奖励，用PPO或GRPO更新\scite{nlp-execution-jin2025searchr1}。

与单轮推理相比，多轮强化学习新增了两个难点。其一是信用分配：一条十几步的轨迹只有末尾一个奖励，第1步改用编号搜索的正确决定与第5步的错误参数得到同一个优势。一种思路是引入轮次层面的价值估计，ArCHer在轮次层面用离策略的时序差分学习一个评价器，再用它给每一轮的词元级策略梯度提供优势\scite{nlp-execution-zhou2024archer}；另一种是为中间步骤设计过程奖励，原理见第\ref{sec:nlp-reasoning-prm}小节。其二是环境成本：组内比较要求从同一初态采样$G$条轨迹，就需要$G$个能重置到相同状态、互相隔离的沙箱；有副作用的真实服务无法这样使用，用户一方往往也要用模型模拟。奖励判定器还必须读取环境状态而不是轨迹中的完成声明，否则策略会学会只写“已完成”，判定方式见第\ref{sec:nlp-execution-recovery}小节。

\subsection{多智能体分工与协作}
\label{sec:nlp-execution-multi-agent}

单个智能体把规划、执行和检查放在同一个上下文中，有两个不便：上下文随任务增长，后期的步骤要在大量无关历史中找信息；同一个上下文检查自己的结果，容易沿用原来的错误判断，Huang等的实验表明，没有外部反馈时模型的自我纠正往往不能改进推理答案\scite{nlp-execution-huang2024selfcorrect}。多智能体方法把任务分给多个模型实例，每个实例有自己的角色提示与较短的上下文，彼此通过消息协作。主要有两类组织方式。

第一类是角色分工。CAMEL让两个实例分别扮演“用户”和“助手”，用户实例按任务逐步下达指令，助手实例执行并汇报，初始提示固定双方角色，防止对话中角色互换或偏离任务\scite{nlp-execution-li2023camel}。MetaGPT把软件开发的标准作业流程写进协作协议：产品经理、架构师、工程师和测试等角色依次产出需求文档、设计文档、代码和测试，角色之间交换结构化文档而非自由对话，各角色从共享消息池订阅与自己相关的消息\scite{nlp-execution-hong2024metagpt}。结构化消息减少了自然语言转述中的信息丢失，流程也规定了谁在什么时候交付什么。

第二类是辩论与聚合。Du等让$N$个实例独立回答同一问题，随后进行$R$轮辩论：每一轮中，每个实例读入其他实例上一轮的答案与理由，修订自己的答案；实例的答案通常在几轮后趋于一致，最后取多数答案\scite{nlp-execution-du2024debate}。记第$i$个实例在第$r$轮的回答为$y_i^{(r)}$，过程为
\[
 \begin{aligned}
 y_i^{(0)}&\sim\pi(\,\cdot\mid x),\qquad
 y_i^{(r)}\sim\pi\bigl(\,\cdot\mid x,\{y_j^{(r-1)}\}_{j=1}^{N}\bigr),\quad r=1,\ldots,R,\\
 \hat y&=\operatorname{majority}\bigl(y_1^{(R)},\ldots,y_N^{(R)}\bigr).
 \end{aligned}
\]
与式\eqref{eq:nlp-reasoning-majority-vote}的自一致性相比，辩论多了实例之间的信息交换：一个实例可以指出另一个实例的计算错误。代价是生成次数为$N(R+1)$，而且第$r$轮的每个上下文都要读入$N$份回答，上下文长度随$N$增长。实例之间互读也可能让错误传播：一个自信的错误答案说服其他实例，多数表决就会集中到错误上。比较时应在相同预算下进行，例如$N=3$、$R=2$共9次生成，应与9次采样的自一致性对照；Huang等在相同回答数下的比较中，辩论并未优于自一致性。

把分工用于本例，可以设规划者、执行者与审查者三个角色。规划者读用户要求，产出目标记录与子目标；执行者是唯一持有写权限的实例，按子目标调用工具；审查者只有只读权限，独立读取\code{R17}、\code{R18}与通知记录，对照目标逐项判定。消息用固定字段：发送者、接收者、消息类型（计划、执行报告、审查结论）、内容，以及所依据的调用编号。终止条件也写进协议：审查者给出通过结论，或协作轮数达到上限，或执行者报告无法满足的约束。单一写入者避免了两个实例重复提交，审查者不读执行者的完成声明，只读环境，相当于把第\ref{sec:nlp-execution-recovery}小节的独立判定做成了一个角色。

多智能体适合子任务需要不同提示、权限或上下文的情形，例如写代码与审代码、执行与审计；单个上下文就能容纳的短任务，增加的通信轮数与调用费用很难收回。判断收益时，应与“单智能体加自检”或“单智能体多次采样”在相同的调用与词元预算下比较。

\subsection{执行恢复与完成判断}
\label{sec:nlp-execution-recovery}

执行中断的原因各不相同，恢复路径也应分开。参数错误退回生成器修正字段；明确的临时故障按约定次数与退避间隔重试；版本冲突要重新读取、检查最新条件再重新规划；信息不足或存在多个候选对象时交给用户澄清，澄清方式见第\ref{chap:nlp-question-answering-dialogue}章。每条路径都设退出条件，不能把所有异常都接到“再试一次”。图\ref{fig:nlp-execution-environment-loop}把语言目标、控制状态与实际环境分开：观察只收录工具真正返回的内容，模型判断“应该已经改好”仍是待验证的假设。

\begin{figure*}[htbp]
  \centering
  % Adapted from academic-drawing/templates/tikz_template.tex.
% Main path uses south -> north ports; feedback is outside the main column.
\begin{tikzpicture}[
  every node/.style={font=\figurefont\small,align=center,text=BookInk},
  block/.style={book node,font=\figurefont\small,
    text width=43mm,minimum height=11mm},
  side/.style={block,text width=29mm},
  flow/.style={book arrow,rounded corners=1.2mm},
  recovery/.style={flow,dashed,draw=BookAmber}
]
  \node[block,draw=FigureInput,fill=FigureInput!10] (goal) at (0,0)
    {目标及保持条件\\对象、授权、目标版本};
  \node[side,draw=FigureInput,fill=FigureInput!10] (user) at (-61mm,0)
    {用户澄清\\或更正};
  \node[block,draw=FigureModel,fill=FigureModel!10] (plan) at (0,-24mm)
    {判断与计划\\依赖、预算、下一步};
  \node[block] (action) at (0,-48mm)
    {绑定参数并行动\\模式、身份、版本检查};
  \node[block,draw=FigureOutput,fill=FigureOutput!10] (environment)
    at (0,-72mm) {实际环境\\记录及副作用};
  \node[side] (external) at (61mm,-72mm)
    {其他入口\\修改共同状态};
  \node[block] (observation) at (0,-96mm)
    {真实工具观察\\字段、回执或异常};
  \node[block] (state) at (0,-120mm)
    {事件更新状态\\事实、假设、未知操作};
  \node[side,draw=BookAmber,fill=BookAmber!8] (recover) at (61mm,-120mm)
    {先回查旧请求\\恢复或重新读取};
  \node[block,draw=FigureOutput,fill=FigureOutput!10] (check) at (0,-144mm)
    {受控读取与结束核验\\目标及保持条件};
  \node[block,draw=FigureOutput,fill=FigureOutput!10] (done) at (0,-168mm)
    {输出完成证据\\或未完成与未知项};
  \draw[flow] (user.east) -- node[above] {更正} (goal.west);
  \draw[flow] (goal.south) -- (plan.north);
  \draw[flow] (plan.south) -- (action.north);
  \draw[flow] (action.south) -- (environment.north);
  \draw[flow] (external.west) -- (environment.east);
  \draw[flow] (environment.south) -- (observation.north);
  \draw[flow] (observation.south) -- (state.north);
  \draw[flow] (state.south) -- (check.north);
  \draw[flow] (check.south) -- (done.north);
  \draw[flow] (state.west) -- (-35mm,-120mm)
    -- (-35mm,-24mm) -- (plan.west);
  \node[text width=22mm,align=right,anchor=east] at (-39mm,-75mm)
    {更新信息\\重选行动};
  \draw[recovery] (state.east) --
    node[above,text width=17mm] {写入未知} (recover.west);
  \draw[recovery] (recover.north) -- (61mm,-96mm) -- (observation.east);
  \node[text width=29mm,align=center] at (61mm,-151mm)
    {回执描述已发生操作\\回读确认当前状态};
\end{tikzpicture}

  \caption{语言要求到环境结果的执行闭环。实线连接行动与真实观察；左侧反馈更新目标并重选行动，右侧虚线处理未知写入。用户更正影响目标，其他入口的修改影响环境，二者通过不同证据进入状态。}
  \label{fig:nlp-execution-environment-loop}
\end{figure*}
\FloatBarrier

写操作超时需要单独的状态。它可能尚未提交，也可能已提交而回执丢失，客户端看到的都是没有确定响应。控制器先记为“结果未知”，查询原请求的状态，再读取对象。模拟服务提供\code{request\_id}：相同请求号与相同参数的重放只返回原回执，实际写入一次；相同请求号换了参数则被拒绝。这称为\term{幂等键}（idempotency key），它依赖服务端持久保存请求与结果，客户端自己生成一个字符串并不能使重试变得安全。

三条教学轨迹核算了恢复的成本。写后超时：第5次调用已把版本1改为2但抛出超时，第6次查询请求得到已提交回执，第7次读取确认字段，共7次调用、1次写入。写前超时：第6次查询请求为不存在，本地服务同步执行、查询立即一致，在这一契约下可以排除请求仍在途；第7次以同一请求号和参数重试，第8次回读，共8次调用、1次写入。真实异步系统中“暂未查到”未必表示未执行，应继续有限回查或报告未知。第三条轨迹叠加了用户更正：第5次写入18日成功但回执超时，随后用户把个人预约改为19日15:00，并从另一入口把备注改成“带电脑”。状态先把目标更新为第2版，同时保留旧请求的未知状态；第6次查询确认旧请求已提交，第7次读到版本3，第8次检查19日时段，第9次以新请求号、版本3只改时间和房间，第10次读到版本4，确认19日15:00与“带电脑”。共10次调用、2次目标不同的写入；重用旧幂等键提交19日会被拒绝，盲目重发旧请求会恢复过时的18日目标，只把回复改成“已改到19日”则不改变任何记录。

完成判断应由独立的判定器读取环境，把“已完成”翻译成可检查的条件。设$g_j(s)$为第$j$项目标或保持条件在终态$s$上是否成立，终态判定为
\begin{equation}
 J(s)=\prod_{j=1}^{m}\mathbb{I}\bigl[g_j(s)\bigr].
 \label{eq:nlp-execution-completion}
\end{equation}
本例的条件是：\code{R17}的时间与地点为目标值，备注符合最新用户更正，\code{R18}不变，未发送通知。乘积要求全部成立，只检查时间就会放过地点仍旧的部分完成。条件取自任务规格与受控读取，不由被评价的执行器依据自己的回复生成；若过程另有要求，如修改前须获授权，则另查事件顺序，不能从终态反推过程合规。WebArena对信息查询比较答案，对导航与配置任务读取页面或存储状态后检查所需内容，不要求行动序列与参考轨迹一致\scite{nlp-intro-zhou2024webarena}；$\tau$-bench同样把对话结束时的数据库状态与标注的目标状态比较\scite{nlp-execution-yao2024taubench}，$\tau^2$-Bench进一步让用户一方也能改变环境\scite{nlp-intro-barres2025}。

表\ref{tab:nlp-execution-outcomes}摘要本地保存的运行记录。“预算为5”一行值得注意：环境已被正确修改，执行器却没有预算做第6次回读，只能报告“已提交，尚未核验”，独立判定随后确认目标达成。环境结果与执行器当时有权声称的结果是两个字段。

\begin{table}[htbp]
  \centering
  \begin{tabular}{lccc}
    \toprule
    轨迹 & 调用数 & 写入数 & 外部目标判定 \\
    \midrule
    正常执行 & 6 & 1 & 通过 \\
    只改时间、未改地点 & 1 & 1 & 不通过 \\
    读取服务不可用 & 3 & 0 & 不通过 \\
    预算为4，查询后停止 & 4 & 0 & 不通过 \\
    预算为5，未做回读 & 5 & 1 & 通过，未自验 \\
    成功修改同名另一对象 & 1 & 1 & 不通过 \\
    只宣布完成、未调用 & 0 & 0 & 不通过 \\
    \bottomrule
  \end{tabular}
  \caption{模拟执行的结果与成本。所有输入和轨迹均为教学构造；独立判定器的读取不计入执行器的调用预算。目标判定不隐藏环境故障或部分完成，停止原因另行记录。}
  \label{tab:nlp-execution-outcomes}
\end{table}

随机的智能体需要在重置后的环境中重复运行。式\eqref{eq:nlp-eval-passk}的pass@$k$衡量$k$次尝试中至少一次成功，适合可以独立重跑、由测试挑选的程序候选；环境任务通常只交付一次，更关心每次是否都成功。$\tau$-bench为此定义pass\textasciicircum{}$k$：对每个任务运行$n$次，其中成功$c$次，
\begin{equation}
 \text{pass\textasciicircum{}}k=\mathbb{E}_{\text{任务}}\left[\binom{c}{k}\Big/\binom{n}{k}\right],
 \label{eq:nlp-execution-pass-hat-k}
\end{equation}
即随机抽取的$k$次运行全部成功的概率。它随$k$增大而下降，下降越快说明行为越不稳定。例如某任务10次运行成功8次，pass\textasciicircum{}1为0.8，pass\textasciicircum{}2为$\binom{8}{2}/\binom{10}{2}=28/45\approx0.62$。两种指标都假定各次运行相互独立；若环境不能无损重置，第二次尝试已受第一次副作用影响，就不能按独立试验估计。成功比例之外，报告还应给出调用数、时延与重复副作用，并按工具选择、参数绑定、状态更新和终态判定分别统计失败。


\section*{文献说明}
\label{sec:nlp-execution-literature}

以下按节列出正文所依据的原始方法与基准入口。思维链、自一致性、思维树、结果与过程奖励模型等推理方法的文献见第\ref{chap:nlp-reasoning}章；pass@$k$的估计见第\ref{chap:nlp-evaluation}章。正文每个方法至多引用一个原文对照数字，其余数值均为教学构造，不复现论文实验。

\begin{fullwidth}
数学应用题的神经求解：深度神经求解器以序列到序列模型生成表达式并给出Math23K数据集；GTS提出目标分解与子树嵌入的树解码器。GSM8K与MATH是正文所说文字数学题的主要基准。
\par\fullcite{nlp-execution-wang2017dns}
\par\fullcite{nlp-execution-xie2019gts}
\par\fullcite{nlp-intro-cobbe2021}
\par\fullcite{nlp-eval-hendrycks2021math}
\end{fullwidth}

\begin{fullwidth}
程序辅助与工具集成推理：PAL以少样本提示生成程序并单次执行；ToRA给出推理、程序与执行输出交替的轨迹格式，以及模仿学习与输出空间塑形的训练流程。Z3命令依据官方指南。
\par\fullcite{nlp-execution-gao2023pal}
\par\fullcite{nlp-execution-gou2024tora}
\par\fullcite{nlp-execution-z3guide}
\end{fullwidth}

\begin{fullwidth}
\raggedright
形式化定理证明：Lean 4的命题、证明项与证明策略依据官方文档；Wu等研究了大语言模型的自动形式化，Lean Workbook以编译与回译一致性筛选译文，DeepSeek-Prover加入假设拒绝并迭代合成证明数据；GPT-f提出策略语言模型、累计对数概率的最佳优先搜索与可证性函数，HTPS在超图上做受AlphaZero启发的搜索，Polu等研究了形式数学上的专家迭代。miniF2F是奥赛级形式命题的跨系统基准。
\par\fullcite{nlp-execution-lean4}
\par\fullcite{nlp-execution-wu2022autoformalization}
\par\fullcite{nlp-execution-ying2024leanworkbook}
\par\fullcite{nlp-execution-xin2024deepseekprover}
\par\fullcite{nlp-execution-polu2020gptf}
\par\fullcite{nlp-execution-lample2022htps}
\par\fullcite{nlp-execution-polu2023curriculum}
\par\fullcite{nlp-execution-zheng2022minif2f}
\end{fullwidth}

\begin{fullwidth}
\raggedright
代码生成与转换：HumanEval给出以单元测试评价函数级代码生成的协议；StarCoder、DeepSeek-Coder与The Stack说明代码模型的训练数据，中间填充目标见Bavarian等；AlphaCode给出大规模采样、过滤与行为聚类；SynCode给出词元与终结符对齐的语法约束解码；Self-Debugging给出执行反馈与代码解释的修订循环；TransCoder以跨语言预训练、去噪自编码与回译做无监督代码翻译。Python AST接口依据官方文档。
\par\fullcite{nlp-intro-chen2021}
\par\fullcite{nlp-execution-li2023starcoder}
\par\fullcite{nlp-execution-guo2024deepseekcoder}
\par\fullcite{nlp-execution-kocetkov2022stack}
\par\fullcite{nlp-llm-fim2022}
\par\fullcite{nlp-execution-li2022alphacode}
\par\fullcite{nlp-execution-ugare2024syncode}
\par\fullcite{nlp-analysis-chen2024selfdebug}
\par\fullcite{nlp-execution-lachaux2020transcoder}
\par\fullcite{nlp-execution-python-ast}
\end{fullwidth}

\begin{fullwidth}
仓库级程序修复：SWE-bench以真实仓库问题和测试评价补丁；SWE-agent提出智能体–计算机接口；Agentless以分层定位、局部修复和补丁验证三阶段流程完成修复。
\par\fullcite{nlp-intro-jimenez2024}
\par\fullcite{nlp-execution-yang2024sweagent}
\par\fullcite{nlp-execution-xia2024agentless}
\end{fullwidth}

\begin{fullwidth}
\raggedright
工具学习与智能体：Toolformer以执行前后的损失差自监督地筛选工具调用；ToolLLM构造ToolBench并以决策树搜索生成调用轨迹，Gorilla以检索到的接口文档为条件微调；ReAct交替判断、行动与观察，Plan-and-Solve与ReWOO分离计划与执行，LLMCompiler按依赖图并行调用；Generative Agents给出按新近度、重要性与相关性检索记忆的打分，MemGPT分层管理上下文与外部记忆，Reflexion以文字反思跨回合积累经验，Huang等讨论了无外部反馈时自我修正的局限；FireAct与AgentTuning以交互轨迹微调智能体，ArCHer与Search-R1在多轮交互上做强化学习；CAMEL、MetaGPT与多智能体辩论代表角色分工与相互检查。$\tau$-Bench、$\tau^2$-Bench与WebArena按环境结果评价智能体。
\par\fullcite{nlp-execution-schick2023toolformer}
\par\fullcite{nlp-execution-qin2024toolllm}
\par\fullcite{nlp-execution-patil2024gorilla}
\par\fullcite{yao2023react}
\par\fullcite{nlp-execution-wang2023planandsolve}
\par\fullcite{nlp-execution-xu2023rewoo}
\par\fullcite{nlp-execution-kim2024llmcompiler}
\par\fullcite{nlp-execution-park2023generative}
\par\fullcite{nlp-execution-packer2023memgpt}
\par\fullcite{nlp-execution-shinn2023reflexion}
\par\fullcite{nlp-execution-huang2024selfcorrect}
\par\fullcite{nlp-execution-chen2023fireact}
\par\fullcite{nlp-execution-zeng2024agenttuning}
\par\fullcite{nlp-execution-zhou2024archer}
\par\fullcite{nlp-execution-jin2025searchr1}
\par\fullcite{nlp-execution-li2023camel}
\par\fullcite{nlp-execution-hong2024metagpt}
\par\fullcite{nlp-execution-du2024debate}
\par\fullcite{nlp-execution-yao2024taubench}
\par\fullcite{nlp-intro-barres2025}
\par\fullcite{nlp-intro-zhou2024webarena}
\end{fullwidth}

\section{本章小结}
\label{sec:nlp-execution-summary}

章首的八项挑战，分别由本章的不同机制承担。语言条件到形式对象的映射由表达式树、约束和形式命题完成：GTS用目标分解与子树嵌入生成合法的表达式树，CSP与SMT把多变量条件交给求解器判定，自动形式化把自然语言命题译成Lean命题并以编译、回译与假设拒绝筛选。可靠计算由执行器承担：PAL把算术交给解释器，ToRA让推理依据真实的执行输出继续。答案与过程的验证按可靠性分层：规则回代独立重算，结果与过程奖励模型在不可规则化的候选间选择并定位错误步骤，证明助手则在命题形式化正确的前提下给出不依赖评分器的保证，证明本身由策略语言模型配合最佳优先搜索、价值估计和专家迭代去寻找。

代码任务把程序变成交付物。代码语言模型的训练数据与中间填充目标决定了基础能力，大规模采样加测试过滤与行为聚类以计算换正确率，语法约束保证可解析，执行反馈让模型修订而不只是重采样；仓库修复的关键在于定位与验证，智能体–计算机接口与分层定位流程代表了两种分配自主决策的方式，多候选补丁还需要按目标测试与回归测试选择。工具任务进一步引入外部状态：调用能力来自工具描述、自监督与合成指令数据，参数要绑定到可靠来源，计划、ReAct交替与依赖图并行分别权衡检查便利、灵活性和等待时间，结构化状态与记忆保存仍影响任务的事实，轨迹微调与多轮强化学习直接训练多步行为，多智能体分工以通信成本换取相互检查。

这些机制都没有单独保证任务完成。形式化正确是计算与证明有意义的前提，测试通过只在测试覆盖的输入范围内说明行为，完成声明则必须对应正确的对象、最新的目标与真实的回读。回到章首“预约是否真的改好”的问题：应当能指出改了哪条记录、依据哪一版目标、保留了哪些字段、读回了什么结果；写入超时而尚未回查时，正确的交付是说明结果未知并给出后续核验。第四部分转向语音任务，输入输出增加了声学信号与时间约束，但取得证据、保存过程与验收目标的区分仍然适用。
